\section{Conventions and group varieties}

\begin{convention}[Standing algebraic-geometric language]
\label{conv:milne-standing-geometry}

\dcref{MI:ch1:1}
A variety over $k$ is a geometrically reduced separated scheme of finite type,
with nonclosed points omitted from its point set.  Base change to a field $K$
is written $V_K$, and $V(K)$ denotes the $K$-coordinate points.  All rings
are commutative with identity and all homomorphisms preserve identity.  The
field $k$ is fixed unless a section explicitly changes it.
\end{convention}

\begin{definition}[Group variety and abelian variety]
\label{def:milne-group-variety}

\dcref{MI:ch1:1}
\lean{MilneLib.GroupVariety.pointsFunctor, MilneLib.GroupVariety.pointsFunctor_representable,
MilneLib.GroupVariety.comp_mulRight_hom, MilneLib.GroupVariety.comp_mulRight_inv,
MilneLib.GroupVariety.pointTranslation, MilneLib.GroupVariety.pointTranslation_self,
MilneLib.GroupVariety.pointTranslation_symm, MilneLib.GroupVariety.pointTranslation_trans,
MilneLib.GroupVariety.comp_pointTranslation_hom, MilneLib.GroupVariety.pointTranslationIso,
MilneLib.GroupVariety.pointTranslationIso_hom, MilneLib.GroupVariety.pointTranslationIso_inv,
MilneLib.GroupVariety.pointTranslationIso_self, MilneLib.GroupVariety.pointTranslationIso_symm,
MilneLib.GroupVariety.pointTranslationIso_trans,
MilneLib.GroupVariety.pointTranslationIso_hom_comp, MilneLib.IsAbelianVariety,
MilneLib.isCommMonObj_of_isAbelianVariety,
MilneLib.smooth_of_isAbelianVariety}
A group variety is a variety equipped with regular multiplication, inverse,
and identity maps satisfying the group diagrams after geometric base change.
Translations are isomorphisms.  A complete connected group variety is an
abelian variety; its group law is written additively.
\end{definition}

\begin{theorem}[Rigidity theorem]
\label{thm:milne-rigidity}

\dcref{MI:ch1:1.1}
\lean{MilneLib.GroupVariety.rigidity_constant_of_two_axes}
\lean{MilneLib.GroupVariety.rigidity_constant_of_two_axes_arbitraryField}
\leanok
Consider a regular map $\alpha:V\times W\to U$, and assume that $V$ is
complete and that $V\times W$ is geometrically irreducible.  If there are
points $u_0\in U(k)$, $v_0\in V(k)$, and $w_0\in W(k)$ such that
\[
 \alpha(V\times\{w_0\})=\{u_0\}=\alpha(\{v_0\}\times W),
\]
then $\alpha(V\times W)=\{u_0\}$.
\end{theorem}

\begin{proof}

\leanok

Since the hypotheses continue to hold after extending scalars from $k$ to
$k^{\rm al}$, we can assume that $k$ is algebraically closed.  Note that $V$
is connected, because otherwise $V\times_k W$ would not be connected, much
less irreducible.  We use the following two facts: if $V$ is complete, then
the projection $q:V\times_kW\to W$ is closed; and if $V$ is complete and
connected, every regular map from $V$ into an affine variety is constant.

Let $U_0$ be an open affine neighbourhood of $u_0$.  Because $q$ is closed,
\[
 Z=q\bigl(\alpha^{-1}(U\setminus U_0)\bigr)
\]
is closed in $W$.  A point $w$ lies outside $Z$ if and only if
$\alpha(V\times\{w\})\subset U_0$.  In particular $w_0\notin Z$, so
$W\setminus Z$ is nonempty.  For $w\in W\setminus Z$, the complete variety
$V\times\{w\}\simeq V$ maps into the affine variety $U_0$, and hence its image
is a point.  This point is $\alpha(v_0,w)=u_0$, so $\alpha$ is constant on the
nonempty open subset $V\times(W\setminus Z)$.  Since $V\times W$ is
irreducible, this subset is dense, and since $U$ is separated, $\alpha$ agrees
with the constant map with value $u_0$ everywhere.

\end{proof}


\begin{corollary}[Morphisms of abelian varieties]
\label{cor:milne-morphism-translation-homomorphism}

\dcref{MI:ch1:1.2}
\lean{MilneLib.GroupVariety.exists_hom_comp_pointTranslation_of_isAbelianVariety_arbitraryField}
\uses{thm:milne-rigidity}
Every regular map $\alpha:A\to B$ of abelian varieties is the composite of a
homomorphism with a translation.
\end{corollary}

\begin{proof}
The regular map $\alpha$ sends the $k$-rational point $0$ of $A$ to a
$k$-rational point $b$ of $B$.  After composing $\alpha$ with translation by
$-b$, we may assume that $\alpha(0)=0$.  Consider
\[
 \phi:A\times A\longrightarrow B,\qquad
 \phi(a,a')=\alpha(a+a')-\alpha(a)-\alpha(a').
\]
This is the difference of the two regular maps obtained from the addition maps
on $A\times A$ and $B\times B$.  We have
$\phi(A\times\{0\})=\{0\}=\phi(\{0\}\times A)$, so the Rigidity Theorem
gives $\phi=0$.  Thus $\alpha$ is a homomorphism, and undoing the translation
proves the assertion.
\end{proof}

\begin{corollary}[Commutativity of an abelian variety]
\label{cor:milne-abelian-commutativity}

\dcref{MI:ch1:1.4}
\lean{MilneLib.GroupVariety.isCommMonObj_of_isAbelianVariety_via_rigidity}
\uses{cor:milne-morphism-translation-homomorphism,thm:milne-rigidity}
The group law on an abelian variety is commutative.
\end{corollary}

\begin{proof}
Commutative groups are distinguished among all groups by the fact that the map
taking an element to its inverse is a homomorphism.  Since the negative map
$a\mapsto-a$ on $A$ takes the zero element to itself, the preceding corollary
shows that it is a homomorphism.  Hence the group law is commutative.
\end{proof}

\begin{corollary}[Product decomposition of maps into an abelian variety]
\label{cor:milne-abelian-product-decomposition}

\dcref{MI:ch1:1.5}
\lean{MilneLib.GroupVariety.existsUnique_hom_additive_decomp_of_rigidity}
\uses{cor:milne-morphism-translation-homomorphism,thm:milne-rigidity}
Let $V$ and $W$ be complete varieties over $k$ with $k$-rational points
$v_0$ and $w_0$, and let $p$ and $q$ be the projection maps
$V\times W\to V$ and $V\times W\to W$.  Let $A$ be an abelian variety.  Then
a morphism $h:V\times W\to A$ such that $h(v_0,w_0)=0$ can be written
uniquely as
\[
 h=f\circ p+g\circ q,
\]
with morphisms $f:V\to A$ and $g:W\to A$ satisfying
$f(v_0)=0$ and $g(w_0)=0$.
\end{corollary}

\begin{proof}
Set
\[
 f=h|_{V\times\{w_0\}},\qquad
 g=h|_{\{v_0\}\times W},
\]
identifying the two axes with $V$ and $W$.  On points, the difference
$\delta=h-(f\circ p+g\circ q)$ is
\[
 \delta(v,w)=h(v,w)-h(v,w_0)-h(v_0,w).
\]
It vanishes on both coordinate axes, so the Rigidity Theorem gives
$\delta=0$.  The restrictions to the two axes also show that any such
decomposition has the same $f$ and $g$, proving uniqueness.
\end{proof}

\section{Complex uniformization and Riemann forms}

\begin{proposition}[Exponential uniformization]
\label{prop:milne-complex-exponential}

\dcref{MI:ch1:2.1}
\uses{def:milne-group-variety}
Let $A$ be an abelian variety of dimension $g$ over $\C$.
\begin{enumerate}
\item There is a unique homomorphism of complex manifolds
\[
 \exp:T_0(A(\C))\longrightarrow A(\C)
\]
such that, for each $v\in T_0(A(\C))$, the map
$z\mapsto\exp(zv)$ is the one-parameter subgroup corresponding to $v$.
The differential of $\exp$ at $0$ is the identity.
\item The map $\exp$ is surjective, and its kernel is a full lattice in the
complex vector space $T_0(A(\C))$.
\end{enumerate}
\end{proposition}

\begin{proof}
The first assertion is the standard construction of the exponential of a
complex Lie group.  It remains to prove (2).  The image $H$ of $\exp$ is a
subgroup of $A(\C)$.  Since $d\exp_0$ is an isomorphism, the inverse function
theorem shows that $\exp$ is a local isomorphism at $0$, so $H$ contains an
open neighbourhood $U$ of $0$.  For $a\in H$, the translate $a+U$ is an open
neighbourhood of $a$ in $H$; hence $H$ is open.  Its complement is a union of
cosets of $H$, so $H$ is also closed.  The complex manifold $A(\C)$ is
connected, and therefore $H=A(\C)$.

Put $V=T_0(A(\C))$, viewed as a real vector space of dimension $2g$.  The
local injectivity of $\exp$ at $0$ gives an open neighbourhood $U$ with
$U\cap\ker(\exp)=\{0\}$, so the kernel is discrete.  A discrete subgroup of a
finite-dimensional real vector space is a lattice.  Finally
$V/\ker(\exp)\simeq A(\C)$ is compact, so the lattice has rank $2g$.
\end{proof}

\begin{theorem}[Cohomology of a real torus]
\label{thm:milne-torus-cohomology}

\dcref{MI:ch1:2.3}
Let $X=V/L$ be a real torus.  There are canonical isomorphisms
\[
 \bigwedge^r H^1(X,\Z)\;\xrightarrow{\sim}\;H^r(X,\Z)
 \;\xrightarrow{\sim}\;\Hom_\Z\!\left(\bigwedge^r L,\Z\right).
\]
\end{theorem}

\begin{proof}
The cup product defines a map
$\bigwedge^rH^1(X,\Z)\to H^r(X,\Z)$.  The Kunneth formula shows that this
map is an isomorphism for products whenever it is an isomorphism for the two
factors; it is plainly an isomorphism for $S^1$, and hence for
$X\simeq(S^1)^n$.  The universal cover $V\to X$ has group of covering
transformations $L$, so
$H^1(X,\Z)\simeq\Hom(\pi_1(X),\Z)=\Hom(L,\Z)$.  The determinant pairing
between $\bigwedge^r\Hom(L,\Z)$ and $\bigwedge^rL$ identifies the former
with the dual of the latter, giving the second isomorphism.
\end{proof}

\begin{definition}[Riemann form and polarizable torus]
\label{def:milne-riemann-form}

\dcref{MI:ch1:2}
An integral alternating form $E$ on a lattice $L\subset V$ is a Riemann form
if its real extension satisfies $E(iv,iw)=E(v,w)$ and the associated Hermitian
form $H(v,w)=E(iv,w)+iE(v,w)$ is positive definite.  A complex torus admitting
such a form is polarizable.  Its dual torus is $V^\vee/L^\vee$, where
$V^\vee$ is the conjugate-linear dual and $L^\vee$ consists of functionals
integral on $L$.
\end{definition}

\begin{lemma}[Hermitian and alternating forms]
\label{lem:milne-riemann-hermitian}

\dcref{MI:ch1:2.4}
Let $V$ be a complex vector space.  There is a one-to-one correspondence
between Hermitian forms $H$ on $V$ and real-valued skew-symmetric forms $E$ on
$V$ satisfying $E(iv,iw)=E(v,w)$, namely
\[
 E(v,w)=\Im(H(v,w)),\qquad
 H(v,w)=E(iv,w)+iE(v,w).
\]
\end{lemma}

\begin{proof}
If $H$ is Hermitian, then $E=\Im H$ is real-valued and skew-symmetric, and
Hermitian symmetry gives $E(iv,iw)=E(v,w)$.  Conversely, if $E$ has these
properties, the displayed formula is additive in the first variable,
conjugate-additive in the second, and satisfies
$H(w,v)=\overline{H(v,w)}$; hence it is Hermitian.  Taking imaginary parts and
substituting $iv$ recovers $E$, so the two constructions are inverse.
\end{proof}

\begin{theorem}[Analytic characterization of abelian varieties]
\label{thm:milne-polarizable-tori}

\dcref{MI:ch1:2.8}
\uses{prop:milne-complex-exponential,thm:milne-torus-cohomology,def:milne-riemann-form}
A complex torus $X$ is of the form $A(\C)$ for a complex abelian variety $A$
if and only if it is polarizable.
\end{theorem}

\begin{proof}
For the forward implication, choose an embedding $A\hookrightarrow\PP^n$
with $n$ minimal and a hyperplane $H$ not containing the tangent space at any
point of $A(\C)$.  The smooth hyperplane section $A\cap H$ defines a class in
\[
 H_{2g-2}(A,\Z)\simeq H^2(A,\Z)
 \simeq\Hom(\bigwedge^2L,\Z),
\]
and the Riemann bilinear relations show that the resulting alternating form
on the period lattice $L$ is a Riemann form.  Conversely, a Riemann form
gives theta functions on $V$ with the required automorphy factors; a suitable
finite collection of these functions separates points and tangent directions
and embeds $V/L$ in projective space.  (Milne labels this reverse direction
as a brief sketch and refers to the theta-function construction.)
\end{proof}

\begin{theorem}[Equivalence with polarizable tori]
\label{thm:milne-complex-category-equivalence}

\dcref{MI:ch1:2.9}
\uses{thm:milne-polarizable-tori}
The functor $A\mapsto A(\C)$ is an equivalence from the category of complex
abelian varieties to the category of polarizable complex tori.  In particular,
it identifies
\[
 \Hom(A,B)=\Hom(A(\C),B(\C)),
\]
where a homomorphism of tori is induced by a complex-linear map carrying the
period lattice into the target lattice.  For $X=V/L$ one has
$X_m=m^{-1}L/L$ and the canonical Weil pairing with the $m$-torsion of the
dual torus; a Riemann form induces the corresponding polarization.
\end{theorem}

\begin{proof}
The preceding theorem gives essential surjectivity.  A holomorphic homomorphism
of complex tori fixing $0$ lifts uniquely to a holomorphic map of universal
covers fixing $0$; its differential is constant, so the lift is complex-linear,
and it carries the source lattice into the target lattice.  Conversely every
such linear map descends.  This proves fullness and faithfulness, while the
descriptions of torsion, duals, and polarizations follow from the quotient and
Riemann-form definitions.
\end{proof}

\section{Rational maps into abelian varieties}

\begin{definition}[Rational and birational maps]
\label{def:milne-rational-map}

\dcref{MI:ch1:3}
A rational map $\phi:V\dashrightarrow W$ is an equivalence class of pairs
$(U,\phi_U)$, where $U$ is a dense open subset of $V$ and
$\phi_U:U\to W$ is regular; two pairs are equivalent when the maps agree on
the intersection of their domains.  The rational map is defined at a point
when one representative is defined there, and its maximal domain of definition
is the union of all such $U$.  A rational map is dominating when the image of
a representative is dense; it then induces an inclusion
$k(W)\hookrightarrow k(V)$ and is birational when this inclusion is an
isomorphism.
\end{definition}

\begin{theorem}[Extension across codimension one]
\label{thm:milne-rational-extension-complete}

\dcref{MI:ch1:3.1}
\uses{def:milne-rational-map}
A rational map $\phi:V\dashrightarrow W$ from a normal variety $V$ to a
complete variety $W$ is defined on an open subset $U$ of $V$ whose complement
$V\setminus U$ has codimension at least two.
\end{theorem}

\begin{proof}
Assume first that $V$ is a curve.  Let $U$ be the domain of a representative
of $\phi$, let $U'$ be its graph in $V\times W$, and let $Z$ be the closure of
$U'$.  The image of $Z$ in $V$ is closed because $W$ is complete, and contains
$U$, hence is all of $V$.  The map $Z\to V$ is a surjective map of complete
nonsingular curves which is an isomorphism over a nonempty open subset; it is
therefore an isomorphism, and the other projection gives the extension.

In general, suppose that a codimension-one point of $V$ lies outside the
maximal domain $U$.  Its local ring is a discrete valuation ring because $V$
is normal.  The map on the generic point extends over this local ring by the
valuative criterion of properness, producing a representative defined on an
open set meeting that codimension-one point, a contradiction.  Thus
$V\setminus U$ has codimension at least two.
\end{proof}

\begin{lemma}[Indeterminacy of group-valued maps]
\label{lem:milne-group-rational-indeterminacy}

\dcref{MI:ch1:3.3}
\uses{def:milne-rational-map}
Let $\phi:V\dashrightarrow G$ be a rational map from a nonsingular variety to
a group variety.  Then either $\phi$ is defined on all of $V$, or the points
where it is not defined form a closed subset of pure codimension one in $V$.
\end{lemma}

\begin{proof}
Define the rational map
\[
 \Phi:V\times V\dashrightarrow G,\qquad
 (x,y)\longmapsto\phi(x)\phi(y)^{-1}.
\]
It is defined at $(x,x)$ if $\phi$ is defined at $x$, and then
$\Phi(x,x)=e$.  Conversely, if $\Phi$ is defined at $(x,x)$, it is defined on
an open neighbourhood of that point.  After shrinking an open set $U$ if
necessary, $\phi$ is defined on $U$ and
\[
 \phi(x)=\Phi(x,u)\phi(u)\qquad(u\in U),
\]
so $\phi$ is defined at $x$.  Thus $\phi$ is defined at $x$ exactly when
$\Phi$ is defined at $(x,x)$.

Let $f$ be a nonzero rational function in the image of $\Phi^*$.  Write
$\operatorname{div}(f)=\operatorname{div}(f)_0-\operatorname{div}(f)_\infty$.
Since $V\times V$ is nonsingular, the local ring at $(x,x)$ consists of the
rational functions whose pole divisor does not contain $(x,x)$.  If $\phi$ is
not defined at $x$, some such $f$ has $(x,x)$ in its pole divisor.  The
intersection of that divisor with the diagonal is either empty or pure of
codimension one in the diagonal.  It contains $x$, and so the indeterminacy
locus contains a codimension-one component through $x$.  This proves the
assertion.
\end{proof}

\begin{theorem}[Rational maps to abelian varieties extend]
\label{thm:milne-rational-to-abelian-regular}

\dcref{MI:ch1:3.2}
\uses{thm:milne-rational-extension-complete,lem:milne-group-rational-indeterminacy}
Every rational map from a nonsingular variety to an abelian variety is regular
on the whole source.
\end{theorem}

\begin{proof}
By the extension theorem the indeterminacy locus has codimension at least two,
whereas the lemma says that a nonempty indeterminacy locus is pure of
codimension one.  Therefore it is empty.
\end{proof}

\begin{theorem}[Two-variable rigidity without completeness]
\label{thm:milne-rational-product-rigidity}

\dcref{MI:ch1:3.4}
\uses{thm:milne-rigidity,thm:milne-rational-to-abelian-regular}
Let $\alpha:V\times W\to A$ be a morphism from a product of nonsingular
varieties into an abelian variety, and assume that $V\times W$ is geometrically
irreducible.  If
\[
 \alpha(V\times\{w_0\})=\{a_0\}=\alpha(\{v_0\}\times W)
\]
for $a_0\in A(k)$, $v_0\in V(k)$, and $w_0\in W(k)$, then
$\alpha(V\times W)=\{a_0\}$.
\end{theorem}

\begin{proof}
We may assume that $k$ is algebraically closed.  If $V$ is a curve, embed it
in a complete nonsingular curve $C$.  By Theorem 3.2, $\alpha$ extends to a
morphism $\bar\alpha:C\times W\to A$, and the Rigidity Theorem shows that
$\bar\alpha$ is constant.  In general, let $C$ be an irreducible curve in $V$
through $v_0$ and nonsingular there, and let $C'\to C$ be its normalization.
The preceding argument shows that $\alpha$ is constant on $C\times W$.  The
union of such curves is dense by the next lemma, so $\alpha$ is constant.
\end{proof}

\begin{lemma}[Compactification and density of curves]
\label{lem:milne-curves-through-smooth-point}

\dcref{MI:ch1:3.5}
\begin{enumerate}
\item Every nonsingular curve $V$ can be realized as an open subset of a
complete nonsingular curve $C$.
\item If $C$ is a curve, there is a nonsingular curve $C'$ and a regular map
$C'\to C$ which is an isomorphism over the nonsingular points of $C$.
\item If $V$ is irreducible over an algebraically closed field and $P$ is a
nonsingular point of $V$, then the union of the irreducible curves passing
through $P$ and nonsingular at $P$ is dense in $V$.
\end{enumerate}
\end{lemma}

\begin{proof}
For (1), take the nonsingular complete curve with function field $k(V)$,
constructed from the discrete valuation rings of $k(V)$ containing $k$.  For
(2), take the normalization of $C$.  For (3), argue by induction on
$\dim V$.  It is enough to show that the union of the irreducible
codimension-one subvarieties through $P$ and nonsingular at $P$ is dense.  Put
$V$ affine and embedded in affine space.  If $H$ is a hyperplane through $P$
which does not contain $T_PV$, the component $V_H$ through $P$ is nonsingular
at $P$.  If a closed set $Z$ contains every $V_H$, then the tangent-cone
inclusions
\[
 T_PV\cap H=T_PV_H\subset C_P(V_H)\subset C_P(Z)\subset C_P(V)=T_PV
\]
for all such $H$ imply $C_P(Z)=T_PV$.  Hence
$\dim Z=\dim C_P(Z)=\dim V$, so $Z=V$.
\end{proof}

\begin{corollary}[Group-valued rational maps]
\label{cor:milne-group-to-abelian}

\dcref{MI:ch1:3.7}
\uses{thm:milne-rational-to-abelian-regular,cor:milne-morphism-translation-homomorphism}
Every rational map $\alpha:G\dashrightarrow A$ from a group variety to an
abelian variety is the composite of a homomorphism $G\to A$ with a translation.
\end{corollary}

\begin{proof}
Theorem 3.2 shows that $\alpha$ is a regular map.  The proof is then the same
as that of Corollary 1.2.
\end{proof}

\begin{theorem}[Birational rigidity of abelian varieties]
\label{thm:milne-birational-abelian-isomorphism}

\dcref{MI:ch1:3.8}
\uses{thm:milne-rational-to-abelian-regular,cor:milne-morphism-translation-homomorphism}
If two abelian varieties are birationally equivalent, then they are isomorphic
as abelian varieties.
\end{theorem}

\begin{proof}
Let $A$ and $B$ be the abelian varieties.  A birational map
$\phi:A\dashrightarrow B$ and its inverse extend to regular maps by Theorem
3.2.  Their composites are the identity on dense open subsets, hence on all
of $A$ and $B$.  Thus $\phi$ is an isomorphism of varieties.  After composing
with a translation it sends $0$ to $0$, and Corollary 1.2 then shows that it
preserves the group structure.
\end{proof}

\begin{proposition}[No rational curves]
\label{prop:milne-no-rational-curves}

\dcref{MI:ch1:3.9}
\uses{thm:milne-rational-to-abelian-regular,cor:milne-morphism-translation-homomorphism}
Every rational map $\A^1\dashrightarrow A$ or $\PP^1\dashrightarrow A$ to an
abelian variety is constant.
\end{proposition}

\begin{proof}
By Theorem 3.2 the map extends to $\A^1$.  After composing with a translation,
assume $\alpha(0)=0$.  Corollary 1.2 gives
$\alpha(x+y)=\alpha(x)+\alpha(y)$.  The multiplicative group
$\A^1\setminus\{0\}$ is also a group variety, so the same corollary gives
$\alpha(xy)=\alpha(x)+\alpha(y)+c$ for a constant $c$.  These two identities
are incompatible for a nonconstant map, so $\alpha$ is constant.  Restriction
to $\A^1\subset\PP^1$ and completeness give the projective-line case.
\end{proof}

\begin{proposition}[Unirational varieties have no nonconstant maps to an AV]
\label{prop:milne-unirational-to-abelian}

\dcref{MI:ch1:3.10}
\uses{prop:milne-no-rational-curves,cor:milne-abelian-commutativity}
Every rational map $\alpha:V\dashrightarrow A$ from a unirational variety to
an abelian variety is constant.
\end{proposition}

\begin{proof}
We may assume $k$ algebraically closed.  Choose a dominant rational map
$\A^n\dashrightarrow V$.  Its composite with $\alpha$ extends to a morphism
$\beta:(\PP^1)^n\to A$.  Induction using Corollary 1.5 gives morphisms
$\beta_i:\PP^1\to A$ such that
$\beta(x_1,\ldots,x_n)=\sum_i\beta_i(x_i)$.  Each $\beta_i$ is constant by
the preceding proposition, so the composite, and hence $\alpha$, is constant.
\end{proof}

\section{Cohomology and base change}

\begin{theorem}[Coherence of higher direct images]
\label{thm:milne-cohomology-coherence}

\dcref{MI:ch1:4.1}
If $f:V\to T$ is a proper regular map and $\mathcal F$ is a coherent
$\mathcal O_V$-module, then the higher direct image sheaves
$R^if_*\mathcal F$ are coherent $\mathcal O_T$-modules for all $i\ge0$.
\end{theorem}

\begin{proof}[Source proof]
When $f$ is projective, this is proved in Hartshorne 1977, III 8.8.  Chow's
lemma (ibid., II Ex. 4.10) allows one to extend the result to the general case
(EGA III 3.2.1).
\end{proof}

\begin{theorem}[Cohomology in a flat family]
\label{thm:milne-cohomology-base-change}

\dcref{MI:ch1:4.2}
\uses{thm:milne-cohomology-coherence}
Let $f:V\to T$ be a proper flat regular map and let $\mathcal F$ be a locally
free $\mathcal O_V$-module of finite rank.  For $t\in T$, write $V_t$ for the
fibre and $\mathcal F_t$ for the inverse image of $\mathcal F$ on $V_t$.
\begin{enumerate}
\item Formation of the higher direct images of $\mathcal F$ commutes with flat
base change.  In particular, if $T=\operatorname{Spec}(R)$ and $R'$ is a
flat $R$-algebra, then
\[
 H^r(V',\mathcal F')=H^r(V,\mathcal F)\otimes_RR'.
\]
\item $t\mapsto\chi(\mathcal F_t)=\sum_r(-1)^r\dim_{k(t)}H^r(V_t,\mathcal F_t)$
is locally constant.
\item For each $r$, $t\mapsto\dim_{k(t)}H^r(V_t,\mathcal F_t)$ is upper
semicontinuous.
\item If $T$ is integral and this function is constant, then $R^rf_*\mathcal F$
is locally free and the natural maps
\[
 R^rf_*\mathcal F\otimes_{\mathcal O_T}k(t)\longrightarrow
 H^r(V_t,\mathcal F_t)
\]
are isomorphisms.
\item If $H^1(V_t,\mathcal F_t)=0$ for all $t$, then $R^1f_*\mathcal F=0$,
$f_*\mathcal F$ is locally free, and its formation commutes with base change.
\end{enumerate}
\end{theorem}

\begin{proof}[Source proof]
(a) The statement is local on the base, and so it suffices to prove it over
$T=\operatorname{Spec}(R)$.  Mumford 1970, \S5, p.46, constructs a complex
$K^\bullet$ of $R$-modules such that, for every $R$-algebra $R'$,
\[
H^r(V',\mathcal F')=H^r(K^\bullet\otimes_RR').
\]
In our case $R'$ is flat over $R$, and therefore
$H^r(K^\bullet\otimes_RR')=H^r(K^\bullet)\otimes_RR'$, which equals
$H^r(V,\mathcal F)\otimes_RR'$.  Parts (b), (c), and (d) are proved in
Mumford 1970, \S5.  For (e), the hypothesis implies that
$R^1f_*\mathcal F=0$ (Hartshorne 1977, III 12.11a), and it follows that
\[
f_*\mathcal F\otimes_{\mathcal O_T}k(t)\longrightarrow
H^0(V_t,\mathcal F_t)
\]
is surjective for all $t$ (ibid., III 12.11b), and hence is an isomorphism.
Now (ibid., III 12.11b), applied with $i=0$, shows that $f_*\mathcal F$ is
locally free.
\end{proof}

\section{The theorem of the cube and seesaw}

\begin{theorem}[Theorem of the cube]
\label{thm:milne-cube}

\dcref{MI:ch1:5.1}
\uses{thm:milne-cohomology-base-change}
Let $U,V,W$ be complete geometrically irreducible varieties over $k$, and let
$u_0\in U(k)$, $v_0\in V(k)$, and $w_0\in W(k)$ be base points.  Then an
invertible sheaf $\mathcal L$ on $U\times V\times W$ is trivial if its
restrictions to
\[
 U\times V\times\{w_0\},\qquad U\times\{v_0\}\times W,\qquad
 \{u_0\}\times V\times W
\]
are all trivial.
\end{theorem}

\begin{proof}[Source proof deferred in Milne]
Milne states the theorem here and explicitly defers its proof to a later
version of the notes.  The subsequent applications use the theorem as an
input; no complete proof is supplied in the reference.
\end{proof}

\begin{corollary}[Cube identities on an abelian variety]
\label{cor:milne-cube-line-bundle-identities}

\dcref{MI:ch1:5.2}
\uses{thm:milne-cube}
Let $p_i:A^3\to A$ be the projection onto the $i$th factor,
$p_{ij}=p_i+p_j$, and $p_{123}=p_1+p_2+p_3$.  For every invertible sheaf $L$
on $A$, the invertible sheaf
\[
p_{123}^*L\otimes p_{12}^*L^{-1}\otimes p_{23}^*L^{-1}
 \otimes p_{13}^*L^{-1}\otimes p_1^*L\otimes p_2^*L\otimes p_3^*L
\]
on $A^3$ is trivial.
\end{corollary}

\begin{proof}
Restrict the displayed sheaf to each of the three coordinate faces through
$0$.  The factors cancel in pairs, so all three restrictions are trivial.  The
Theorem of the Cube therefore makes the displayed sheaf trivial.
\end{proof}

\begin{theorem}[Theorem of the square]
\label{thm:milne-square}

\dcref{MI:ch1:5.5}
\uses{cor:milne-cube-line-bundle-identities}
For every invertible sheaf $L$ on $A$ and points $a,b\in A(k)$,
\[
 t_{a+b}^*L\otimes L\simeq t_a^*L\otimes t_b^*L.
\]
\end{theorem}

\begin{proof}
Apply Corollary~\ref{cor:milne-cube-line-bundle-identities} to the maps
$x\mapsto x$, $x\mapsto a$, and $x\mapsto b$ from $A$ to itself.  The
resulting triviality is
\[
t_{a+b}^*L\otimes t_a^*L^{-1}\otimes t_b^*L^{-1}\otimes L\simeq\mathcal O_A,
\]
which is the displayed formula.  Tensoring with $L^{-1}$ shows that
$a\mapsto t_a^*L\otimes L^{-1}$ is a homomorphism.
\end{proof}

\begin{theorem}[Seesaw principle]
\label{thm:milne-seesaw}

\dcref{MI:ch1:5.16}
\uses{thm:milne-cohomology-base-change}
Let $V$ and $T$ be varieties over $k$ with $V$ complete, and let $L$ be an
invertible sheaf on $V\times T$.  If $L_t$ is trivial for all $t\in T$, then
there exists an invertible sheaf $N$ on $T$ such that
$L\simeq q^*N$, where $q:V\times T\to T$ is the projection.
\end{theorem}

\begin{proof}
By assumption $H^0(V_t,L_t)\simeq k(t)$ for every $t$.  The cohomology and
base-change theorem therefore shows that $N=q_*L$ is invertible.  The canonical
evaluation map $q^*N=q^*q_*L\to L$ restricts on each fibre to
\[
 \mathcal O_{V_t}\otimes_{k(t)}H^0(V_t,\mathcal O_{V_t})\longrightarrow
 \mathcal O_{V_t},
\]
which is an isomorphism.  Nakayama's lemma shows that the evaluation map is an
isomorphism, so $L\simeq q^*N$.
\end{proof}

\begin{corollary}[Seesaw reduction]
\label{cor:milne-seesaw-twist}

\dcref{MI:ch1:5.17}
Let $V$ and $T$ be varieties over $k$ with $V$ complete, and let $L$ and $M$
be invertible sheaves on $V\times T$.  If $L_t\simeq M_t$ for all $t\in T$,
then there exists an invertible sheaf $N$ on $T$ such that
$L\simeq M\otimes q^*N$.
\end{corollary}

\begin{proof}
Apply Theorem~\ref{thm:milne-seesaw} to $L\otimes M^{-1}$.
\end{proof}

\begin{corollary}[Seesaw principle]
\label{cor:milne-seesaw-principle}

\dcref{MI:ch1:5.18}
Under the hypotheses of the preceding corollary, if $L_v\simeq M_v$ for at
least one $v\in V(k)$, then $L\simeq M$.
\end{corollary}

\begin{proof}
Write $L\simeq M\otimes q^*N$ by the preceding corollary.  Pulling back along
$t\mapsto(v,t)$ gives $L_v\simeq M_v\otimes N$.  The assumed isomorphism
$L_v\simeq M_v$ therefore makes $N$ trivial.
\end{proof}

\begin{proposition}[Closed locus of trivial fibres]
\label{prop:milne-seesaw-closed-locus}

\dcref{MI:ch1:5.19}
Let $V$ be a complete variety, and let $L$ be an invertible sheaf on
$V\times T$.  Then
\[
\{t\in T:L_t\text{ is trivial}\}
\]
is closed in $T$.
\end{proposition}

\begin{proof}
The locus is the intersection of $\operatorname{Supp}(q_*L)$ and
$\operatorname{Supp}(q_*L^\vee)$, which are closed by properness.
\end{proof}

\begin{lemma}[Curves through two points]
\label{lem:milne-curve-through-two-points}

\dcref{MI:ch1:5.20}
Let $V$ be an irreducible variety over an algebraically closed field, and let
$P,Q\in V$.  Then there is an irreducible curve on $V$ passing through both
$P$ and $Q$.
\end{lemma}

\begin{proof}[Source proof]
If $V$ is a curve or $P=Q$, there is nothing to prove.  Otherwise, Chow's
lemma gives a projective variety $W$ and a surjective birational morphism
$W\to V$, so we may assume that $V$ is projective.  Blow up $V$ at
$\{P,Q\}$, and choose a general hyperplane section of the resulting
projective variety.  Bertini's theorem makes this section irreducible, while
dimension counting forces it to meet both exceptional divisors.  Its image in
$V$ is the required proper irreducible subvariety; induction on $\dim V$
gives a curve through $P$ and $Q$.
\end{proof}

\section{Projectivity}

\begin{definition}[Very ample linear system]
\label{def:milne-very-ample}

\dcref{MI:ch1:6}
A divisor is very ample when its complete linear system defines a closed
immersion into projective space; it is ample when a positive multiple is very
ample.
\end{definition}

\begin{theorem}[Every abelian variety is projective]
\label{thm:milne-abelian-projective}

\dcref{MI:ch1:6.4}
\uses{thm:milne-square,def:milne-very-ample}
Every abelian variety $A$ is projective.
\end{theorem}

\begin{proof}
First suppose that $k$ is algebraically closed.  We construct prime divisors
$Z_1,\ldots,Z_n$ on $A$ such that
\[
 \bigcap_i Z_i=\{0\},\qquad \bigcap_iT_0(Z_i)=\{0\}
 \subset T_0(A).
\]
Any two points $0$ and $P$ of $A$ lie in an affine open: if $U$ is affine
around $0$, choose $u\in U\cap(U+P)$ and use the translate
$U+P-u$.  Embed this affine open in affine space and choose a hyperplane
through $0$ not through $P$; the closure of its intersection with $A$ is a
prime divisor through $0$ but not $P$.  Repeating and using the descending
chain condition on closed subsets gives the first equality.  Applying the same
argument in an affine neighbourhood of the origin to a tangent vector gives
the second equality.

Let $D=\sum_iZ_i$.  The theorem of the square gives, for arbitrary points
$a_i,b_i\in A$, the linear equivalence
\[
 \sum_i\bigl(Z_{i,a_i}+Z_{i,b_i}+Z_{i,a_i-b_i}\bigr)\sim3D,
\]
where $Z_{i,a}$ denotes the translate by $a$.  Let $a\ne b$.  Choose $Z_1$
not containing $b-a$ and put $a_1=a$.  The choices of $b_1$ for which
$Z_{1,b_1}$ or $Z_{1,a-b_1}$ contains $b$ form proper closed subsets, and the
remaining $a_i,b_i$ can be chosen so that none of the corresponding translates
contains $b$.  The displayed divisor therefore contains $a$ but not $b$.
Thus $3D$ separates points.  The identical construction with a hyperplane
avoiding a prescribed tangent direction shows that $3D$ separates tangent
directions.  The point-and-tangent criterion for a base-point-free linear
system now gives a closed immersion of $A$ into projective space.

For a general field, extend scalars to $k^{\rm al}$ and apply the construction.
Proposition~\ref{prop:milne-6-6}(d) descends the existence of an ample divisor
to $A$, so a positive multiple gives a projective embedding over $k$.
\end{proof}

\section{Isogenies and torsion}

\begin{definition}[Isogeny]
\label{def:milne-isogeny}

\dcref{MI:ch1:7}
\lean{MilneLib.isogenyKernel, MilneLib.isogenyKernelToBase, MilneLib.Isogeny,
MilneLib.Isogeny.surjective, MilneLib.Isogeny.finite_kernel,
MilneLib.Isogeny.kernel_baseChange_isFinite}
The kernel of a homomorphism $\alpha:A\to B$ of abelian varieties is the fibre
over $0$ in the sense of algebraic spaces.  The homomorphism $\alpha$ is an
isogeny if it is surjective and has finite kernel.  Its degree is its degree as
a regular map.  Write $A_n=\ker(n_A)$ for the $n$-torsion group scheme.
\end{definition}

\begin{proposition}[Characterization of isogenies]
\label{prop:milne-isogeny-characterization}

\dcref{MI:ch1:7.1}
\uses{def:milne-isogeny,thm:milne-abelian-projective}
For a homomorphism $\alpha:A\to B$ of abelian varieties, the following are
equivalent:
\begin{enumerate}
\item $\alpha$ is an isogeny;
\item $\dim A=\dim B$ and $\alpha$ is surjective;
\item $\dim A=\dim B$ and $\ker(\alpha)$ is finite;
\item $\alpha$ is finite, flat, and surjective.
\end{enumerate}
\end{proposition}

\begin{proof}
Because $A$ is complete, $\alpha(A)$ is closed in $B$.  Translation by a point
of $\alpha(A)$ identifies every fibre over $\alpha(A)$ with the fibre over
$0$.  The fibre-dimension theorem therefore gives
\[
 \dim\ker(\alpha)=\dim A-\dim\alpha(A).
\]
This proves the equivalence of (a), (b), and (c).  Clearly (d) implies (a).
Conversely, under (a) all fibres have dimension zero, so $\alpha$ is
quasi-finite.  Since $A$ is complete and $B$ is separated, $\alpha$ is proper;
a proper quasi-finite morphism is finite.  Finally $\alpha_*\mathcal O_A$ is
coherent and locally free, so $\alpha$ is flat.
\end{proof}

\begin{theorem}[Degree of multiplication]
\label{thm:milne-multiplication-degree}

\dcref{MI:ch1:7.2}
\uses{def:milne-isogeny}
Let $A$ be an abelian variety of dimension $g$, and let $n>0$.  Then
$n_A:A\to A$ is an isogeny of degree $n^{2g}$.  It is always $\etale$ when
$\operatorname{char}(k)=0$, and when $\operatorname{char}(k)=p>0$ it is $\etale$
if and only if $p$ does not divide $n$.
\end{theorem}

\begin{proof}
From Theorem~\ref{thm:milne-abelian-projective} and Proposition
\ref{prop:milne-6-6}, choose a symmetric very ample invertible sheaf $L$ on
$A$.  Corollary~\ref{cor:milne-cube-line-bundle-identities} gives
$n_A^*L\simeq L^{\otimes n^2}$, which is again very ample.  If
$Z=\ker(n_A)$, then
$(n_A^*L)|_Z\simeq L^{\otimes n^2}|_Z$ is both ample and trivial on the
connected component of $Z$ through $0$.  A trivial invertible sheaf on a
connected complete variety is ample only when that variety is a point; hence
$Z$ is finite.

Write $L=L(D)$ with $D$ a very ample divisor.  Then
\[
 (n_A^*D)^g=\deg(n_A)D^g,\qquad n_A^*D\sim n^2D,
\]
and hence $\deg(n_A)=n^{2g}$, since $D^g\ne0$.  Indeed, the embedding defined
by $D$ admits hyperplanes $H_1,\ldots,H_g$ whose intersections with $A$ are
proper, and
\[
((H_1\cap A)\cdots(H_g\cap A))=\deg(A)\ne0.
\]
Finally, the differential of $n_A$ at $0$ is multiplication by $n$ on
$T_0(A)$.  It is an isomorphism exactly when the characteristic does not
divide $n$; in that case $n_A$ is etale at $0$, hence everywhere by
homogeneity, and otherwise it is not etale.
\end{proof}

\section{Duality and the Picard variety}

\begin{definition}[Dual abelian variety and Poincare bundle]
\label{def:milne-dual-abelian}

\dcref{MI:ch1:8}
For an abelian variety $A$, a pair $(A^\vee,P)$ consisting of a variety
$A^\vee$ and an invertible sheaf $P$ on $A\times A^\vee$ is called the dual
abelian variety and the Poincare sheaf if:
\begin{enumerate}
\item $P|_{A\times\{b\}}\in\Pic^0(A)$ for every $b\in A^\vee$;
\item $P|_{\{0\}\times A^\vee}$ is trivial; and
\item for every variety $T$ and invertible sheaf $L$ on $A\times T$ satisfying
the analogous two conditions, there is a unique regular map $T\to A^\vee$
whose pullback of $P$ is $L$.
\end{enumerate}
Once $A^\vee$ exists, $A^\vee(k)=\Pic^0(A)$.  For an invertible sheaf $L$
on $A$, the theorem of the square gives a homomorphism
\[
 \lambda_L:A\longrightarrow A^\vee,
 \qquad a\longmapsto[t_a^*L\otimes L^{-1}].
\]
\end{definition}

\begin{proposition}[Algebraically trivial line bundles]
\label{prop:milne-pic0-characterizations}

\dcref{MI:ch1:8.4}
\uses{thm:milne-square}
For an invertible sheaf $L$ on $A$, put
\[
 K(L)=\{a\in A:t_a^*L\simeq L\}.
\]
The following conditions are equivalent:
\[
 K(L)=A,\qquad t_a^*L\simeq L\text{ on }A_{k^{\rm al}}
 \text{ for all }a\in A(k^{\rm al}),\qquad
 m^*L\simeq p^*L\otimes q^*L.
\]
\end{proposition}

\begin{proof}
The equivalence of the first two conditions is the definition of $K(L)$.  If
the third condition holds, restrict it to $A\times\{a\}$ to obtain
$t_a^*L\simeq L$.  Conversely, assume that $K(L)=A$.  The line bundle
$m^*L\otimes p^*L^{-1}\otimes q^*L^{-1}$ is trivial on every fibre of
$A\times A\to A$ and on $\{0\}\times A$; the Seesaw Principle therefore
gives its triviality, which is the third condition.
\end{proof}

\begin{proposition}[Ampleness criterion for the dual map]
\label{prop:milne-dual-ampleness-criterion}

\dcref{MI:ch1:8.1}
\uses{def:milne-dual-abelian,thm:milne-seesaw}
Let $L$ be an invertible sheaf on $A$ such that $h^0(A,L)\ne0$.  Then $L$ is
ample if and only if $K(L)$ has dimension zero.
\end{proposition}

\begin{proof}[Source proof partially supplied]
We may assume that $k$ is algebraically closed.  Suppose that $L$ is ample,
and let $B$ be the connected component of $K(L)$ passing through $0$.  It is
an abelian variety (possibly zero), and $L_B=L|_B$ is ample.  Since
$t_b^*L_B\simeq L_B$ for all $b\in B$, Proposition
\ref{prop:milne-pic0-characterizations} gives
\[
m^*L_B\otimes p^*L_B^{-1}\otimes q^*L_B^{-1}\simeq\mathcal O_{B\times B}.
\]
Pulling this back along $b\mapsto(b,b)$ gives
$L_B\otimes(-1)_B^*L_B\simeq\mathcal O_B$.  The left side is ample, so
$B$ has dimension zero.  Hence $K(L)$ has dimension zero.

Milne explicitly notes that the converse is to be added; no converse proof is
present in the reference.
\end{proof}

\begin{definition}[Divisorial correspondence]
\label{def:milne-divisorial-correspondence}
\dcref{MI:ch1:8.0}
An invertible sheaf $L$ on $A\times B$ is a divisorial correspondence between
abelian varieties $A$ and $B$ if its restrictions to
$\{0\}\times B$ and $A\times\{0\}$ are both trivial.  If
$s:A\times B\to B\times A$ switches the factors, then $s^*L$ is a
divisorial correspondence in the opposite direction.
\end{definition}

\begin{theorem}[Divisorial criterion for the dual correspondence]
\label{thm:milne-dual-divisorial-criterion}

\dcref{MI:ch1:8.9}
\uses{def:milne-dual-abelian,prop:milne-dual-ampleness-criterion}
Let $L$ be a divisorial correspondence between abelian varieties $A$ and $B$.
The following are equivalent:
\begin{enumerate}
\item $(B,L)$ is the dual abelian variety of $A$;
\item if $L|_{A\times\{b\}}$ is trivial, then $b=0$;
\item if $L|_{\{a\}\times B}$ is trivial, then $a=0$;
\item $(A,s^*L)$ is the dual abelian variety of $B$, where $s$ switches the
factors.
\end{enumerate}
\end{theorem}

\begin{proof}[Source proof omitted]
Milne calls this result not difficult and refers to Mumford 1970, p.81; no
proof is included in the reference.
\end{proof}

\begin{proposition}[Quotients by finite group actions]
\label{prop:milne-finite-quotient}

\dcref{MI:ch1:8.10}
Let $V$ be an algebraic variety over an algebraically closed field $k$, and let
a finite group $G$ act on $V$ by regular maps.  Assume that every orbit of $G$
is contained in an open affine subset of $V$.  Then there exists a variety
$W=V/G$ and a finite regular map $\pi:V\to W$ such that:
\begin{enumerate}
\item $(W,\pi)$ is the topological quotient of $V$ by $G$;
\item for every open affine $U\subset W$,
$(U,\mathcal O_U)=(\pi^{-1}U,\mathcal O_V)^G$.
\end{enumerate}
The pair is unique up to unique isomorphism, $\pi$ is surjective, and $\pi$
is $\etale$ if $G$ acts freely.
\end{proposition}

\begin{proof}[Source proof omitted]
Milne refers to Mumford 1970, p.66, or Serre 1959, p.57, for the quotient
construction and its properties.
\end{proof}

\begin{proposition}[Surjectivity criterion for the polarization map]
\label{prop:milne-lambda-ample}

\dcref{MI:ch1:8.14}
\uses{def:milne-dual-abelian,thm:milne-abelian-projective}
If $L$ is ample, then $\lambda_L:A\to A^\vee=\Pic^0(A)$ is surjective.
\end{proposition}

\begin{proof}[Source proof omitted]
Milne refers to Mumford 1970, \S8, p.77, or Lang 1959, p.99, for this result.
\end{proof}

\begin{theorem}[Existence of the dual abelian variety]
\label{thm:milne-dual-existence}
\dcref{MI:ch1:8.15}
For every abelian variety $A$, there is a pair $(A^\vee,P)$ satisfying the
universal property in Definition~\ref{def:milne-dual-abelian}; it is unique
up to unique isomorphism.
\end{theorem}

\begin{proof}[Source proof omitted]
Milne constructs $A^\vee$ as the quotient of $A$ by the subgroup scheme
$K(L)$ attached to an ample sheaf $L$ and refers to Mumford, Chapter III, for
the full universal-property proof.
\end{proof}

\section{Dual exactness}

\begin{theorem}[Dual exact sequence]
\label{thm:milne-dual-exact-sequence}

\dcref{MI:ch1:9.1}
\uses{def:milne-dual-abelian,def:milne-isogeny}
If $\alpha:A\to B$ is an isogeny with kernel $N$, then
$\alpha^\vee:B^\vee\to A^\vee$ is an isogeny with kernel $N^\vee$, the
Cartier dual of $N$.  In other words, the exact sequence
\[
0\longrightarrow N\longrightarrow A\xrightarrow{\alpha}B\longrightarrow0
\]
gives rise to the dual exact sequence
\[
0\longrightarrow N^\vee\longrightarrow B^\vee\xrightarrow{\alpha^\vee}
A^\vee\longrightarrow0.
\]
\end{theorem}

\begin{proof}[Source proof omitted]
Milne refers to Mumford, section 15 (algebraically closed case), and Oort for
the general case; the proof is not included in these notes.
\end{proof}

\begin{proposition}[Extensions by the multiplicative group]
\label{prop:milne-pic0-ext}

\dcref{MI:ch1:9.2}
\uses{def:milne-dual-abelian}
Let $L$ be an invertible sheaf on $A$ whose class is in $\Pic^0(A)$, and let
$G(L)$ denote its total space with the zero section removed.  For some
$k$-rational point $e\in G(L)$, the induced map over addition makes $G(L)$ a
commutative group variety with identity $e$, and
\[
0\longrightarrow\Gm\longrightarrow G(L)\longrightarrow A\longrightarrow0
\]
is an exact sequence.
\end{proposition}

\begin{proof}[Source proof omitted]
Milne states that the theorem of the square supplies the compatible fibrewise
multiplication and refers to Serre, VII section 3, for the verification of the
group law and exactness.
\end{proof}

\begin{proposition}[Picard extensions]\label{prop:milne-pic0-ext-classification}
\dcref{MI:ch1:9.3}
The assignment $L\mapsto E(L)$ defines an isomorphism
\[
\Pic^0(A)\xrightarrow{\sim}\Ext^1_k(A,\Gm).
\]
\end{proposition}

\begin{proof}[Source proof omitted]
Milne refers to Serre, VII section 3, for the identification of extensions
with algebraically trivial line bundles.
\end{proof}

\section{Endomorphisms and Tate modules}

\begin{definition}[Tate module and rational endomorphisms]
\label{def:milne-tate-module}

\dcref{MI:ch1:10}
For a prime $\ell\ne\operatorname{char}(k)$,
$T_\ell A=\varprojlim_n A_{\ell^n}(k^{\mathrm{sep}})$ and
$V_\ell A=T_\ell A\otimes_{\Z_\ell}\Q_\ell$.  Write
$\End^0(A)=\End(A)\otimes\Q$.
\end{definition}

\begin{definition}[Simple abelian variety]
\label{def:milne-simple-abelian-variety}

\dcref{MI:ch1:10}
An abelian variety $A$ is simple if there is no abelian variety $B$ with
$0\ne B\ne A$ and $B\subset A$.
\end{definition}

\begin{proposition}[Isogeny decomposition]
\label{prop:milne-simple-isogeny-decomposition}

\dcref{MI:ch1:10.1}
\uses{def:milne-isogeny,def:milne-simple-abelian-variety,prop:milne-lambda-ample}
For any abelian variety $A$, there are simple abelian subvarieties
$A_1,\ldots,A_n\subset A$ such that the map
\[
A_1\times\cdots\times A_n\longrightarrow A,
\qquad (a_1,\ldots,a_n)\longmapsto a_1+\cdots+a_n
\]
is an isogeny.
\end{proposition}

\begin{proof}
Induct on $\dim A$.  It is enough to prove that if $0\ne B\ne A$ is an
abelian subvariety, then there is an abelian subvariety $B'\subset A$ such
that $(b,b')\mapsto b+b'$ is an isogeny $B\times B'\to A$.  Let
$i:B\hookrightarrow A$, choose an ample sheaf $L$ on $A$, and let $B'$ be
the connected component through $0$ of the kernel of
\[
A\xrightarrow{\lambda_L}A^\vee\xrightarrow{i^\vee}B^\vee.
\]
The dimension estimate gives $\dim B'\geq\dim A-\dim B$.  The restriction
of this composite to $B$ is $\lambda_{L|B}$, whose kernel is finite because
$L|B$ is ample.  Hence $B\cap B'$ is finite, and the sum map is an
isogeny.  Induction on $B$ and $B'$ gives the asserted simple factors.
\end{proof}

\begin{proposition}[Tate-module rank]
\label{prop:milne-tate-rank}

\dcref{MI:ch1:10.5}
\uses{def:milne-tate-module,prop:milne-isogeny-characterization}
For a prime $\ell\ne\operatorname{char}(k)$, $T_\ell A$ is a free
$\Z_\ell$-module of rank $2g$, where $g=\dim A$.
\end{proposition}

\begin{proof}
For every $n$, $A_{\ell^n}(k^{\rm sep})$ is a free
$\Z/\ell^n\Z$-module of rank $2g$.  The transition maps are the canonical
quotients, so the inverse limit is a free $\Z_\ell$-module of rank $2g$.
\end{proof}

\begin{lemma}[Torsion groups of constant rank]
\label{lem:milne-torsion-group}

\dcref{MI:ch1:10.3}
Let $Q$ be a torsion abelian group, and write $Q_n$ for the subgroup of
elements killed by $n$.  If $|Q_n|=n^d$ for every positive integer $n$, then
$Q\simeq(\Q/\Z)^d$.
\end{lemma}

\begin{proof}
For each $n$, the hypothesis implies $Q_n\simeq(\Z/n\Z)^d$.  Choose a
cofinal sequence $n_1\mid n_2\mid\cdots$ in which every positive integer
divides some $n_i$, and choose bases compatibly under the inclusions
$Q_{n_i}\subset Q_{n_{i+1}}$.  Their union identifies
$Q=\varinjlim_iQ_{n_i}$ with $(\Q/\Z)^d$.
\end{proof}

\begin{lemma}[The $\ell$-primary form]
\label{lem:milne-ell-primary-torsion}

\dcref{MI:ch1:10.4}
If $Q$ is an $\ell$-primary torsion group and
$|Q_{\ell^n}|=(\ell^n)^d$ for all $n>0$, then
$Q\simeq(\Q_\ell/\Z_\ell)^d$.
\end{lemma}

\begin{proof}
Use the preceding compatible-basis construction on the chain
$\ell^{-n}\Z/\Z\subset\Q_\ell/\Z_\ell$ and pass to the direct limit.
\end{proof}

\begin{lemma}[Injectivity on Tate modules]
\label{lem:milne-tate-injective}

\dcref{MI:ch1:10.6}
\uses{def:milne-tate-module,prop:milne-simple-isogeny-decomposition}
For any prime $\ell\ne\operatorname{char}(k)$, the natural map
\[
\Hom(A,B)\longrightarrow\Hom_{\Z_\ell}(T_\ell A,T_\ell B)
\]
is injective.  In particular, $\Hom(A,B)$ is torsion-free.
\end{lemma}

\begin{proof}
Let $\alpha$ induce the zero map on $T_\ell$.  Then $\alpha$ kills every
prime-to-characteristic torsion point of $A$.  If $A_0\subset A$ is simple,
the kernel of $\alpha|_{A_0}$ is not finite, since it contains all
$A_{0,\ell^n}$; hence $\alpha|_{A_0}=0$.  The sum map from the simple factors
in Proposition~\ref{prop:milne-simple-isogeny-decomposition} is surjective,
so $\alpha=0$.
\end{proof}

\begin{lemma}[Index and determinant]
\label{lem:milne-index-determinant}

\dcref{MI:ch1:10.8}
\lean{MilneLib.LinearMap.natCard_quotient_range_eq_natAbs_det_of_baseChange_bijective}
\leanok
Let $\alpha$ be an endomorphism of a free $\Z$-module $\Lambda$ of finite
rank such that $\alpha\otimes1:\Lambda\otimes\Q\to\Lambda\otimes\Q$ is an
isomorphism.  Then
\[
(\Lambda:\alpha\Lambda)=|\det(\alpha)|.
\]
\end{lemma}

\begin{proof}
\leanok
Choose bases of $\Lambda$ and use Smith normal form for the integral matrix
of $\alpha$.  After multiplying on the left and right by unimodular matrices,
the matrix is diagonal with nonzero entries $d_1,\ldots,d_r$.  The index is
$\prod_i|d_i|$, which is $|\det(\alpha)|$.
\end{proof}

\begin{theorem}[Characteristic polynomial and degree]
\label{thm:milne-endomorphism-characteristic-polynomial}

\dcref{MI:ch1:10.9}
\uses{def:milne-isogeny,thm:milne-square,prop:milne-degree-polynomial-function}
Let $\alpha\in\End(A)$.  There is a unique monic polynomial
$P_\alpha\in\Z[X]$ of degree $2g$ such that
\[
P_\alpha(r)=\deg(\alpha-r_A)
\qquad\text{for all integers }r.
\]
\end{theorem}

\begin{proof}[Proof of Theorem~\ref{thm:milne-endomorphism-characteristic-polynomial}]
Extend the degree to $\End^0(A)$ by
$\deg(\alpha)=n^{-2g}\deg(n\alpha)$ when $n\alpha\in\End(A)$.  For fixed
$\alpha,\beta$, put $D_n=(n\alpha+\beta)^*D$ for a very ample divisor $D$.
The square identity gives
\[
D_{n+2}-2D_{n+1}+D_n\sim D_0,
\]
with $D_0$ independent of $n$.  Intersecting with itself $g$ times shows
$n\mapsto\deg(n\alpha+\beta)$ is a polynomial of degree at most $2g$.
Thus $\deg$ is a homogeneous polynomial function of degree $2g$ on
$\End^0(A)$, and $P_\alpha(X)=\deg(\alpha-X_A)$ is a polynomial of degree
$2g$.  If $D$ is ample and symmetric, the same intersection calculation
shows that $P_\alpha$ is monic with integral coefficients for
$\alpha\in\End(A)$.  A polynomial agreeing with the degree at all integers is
unique.
\end{proof}

\begin{lemma}[Polynomial functions]
\label{lem:milne-polynomial-function}

\dcref{MI:ch1:10.12}
Let $V$ be a vector space over an infinite field $K$, and let $f:V\to K$ be
a function such that, for all $v,w\in V$, the function
$x\mapsto f(xv+w)$ is a polynomial in $x$ with coefficients in $K$.  Then
$f$ is a polynomial function.
\end{lemma}

\begin{proof}
We prove by induction on $n$ that, for every $v_1,\ldots,v_n,w$, the
function $f(x_1v_1+\cdots+x_nv_n+w)$ is a polynomial in the $x_i$.  The case
$n=1$ is the hypothesis.  For the induction step, write the polynomial in
$x_n$ as $\sum_{i=0}^d a_i(x_1,\ldots,x_{n-1})x_n^i$.  Choose distinct
$c_0,\ldots,c_d\in K$ and solve the resulting Vandermonde system for the
$a_i$.  Each $a_i$ is a linear combination of the functions obtained by
putting $x_n=c_j$, which are polynomials by induction.
\end{proof}

\begin{proposition}[Degree as a polynomial function]
\label{prop:milne-degree-polynomial-function}

\dcref{MI:ch1:10.13}
\uses{lem:milne-polynomial-function,thm:milne-square}
The function
\[
\End^0(A)\longrightarrow\Q,\qquad \alpha\longmapsto\deg(\alpha),
\]
is a homogeneous polynomial function of degree $2g$.
\end{proposition}

\begin{proof}
By Lemma~\ref{lem:milne-polynomial-function}, it is enough to show that
$\deg(n\alpha+\beta)$ is a polynomial in $n$ for fixed $\alpha,\beta$.
After clearing denominators, take $\alpha,\beta\in\End(A)$ and let $D$ be a
very ample divisor.  Put $D_n=(n\alpha+\beta)^*D$.  The theorem of the cube
gives
\[
  D_{n+2}-2D_{n+1}+D_n\sim D',
  \qquad D'=(2\alpha)^*D-2\alpha^*D.
\]
Induction gives
\[
 D_n=\frac{n(n-1)}2D'+nD_1-(n-1)D_0,
\]
and intersecting $D_n$ with itself $g$ times shows that
$\deg(n\alpha+\beta)$ is a polynomial of degree at most $2g$.  Homogeneity
follows from $\deg(n\alpha)=n^{2g}\deg(\alpha)$.
\end{proof}

\begin{proposition}[Characteristic polynomial on the Tate module]
\label{prop:milne-endomorphism-tate-characteristic}

\dcref{MI:ch1:10.20}
\uses{thm:milne-endomorphism-characteristic-polynomial,def:milne-tate-module,lem:milne-ell-adic-root-separation,lem:milne-multiplicative-polynomial}
For every prime $\ell\ne\operatorname{char}(k)$, $P_\alpha(X)$ is the
characteristic polynomial of $\alpha$ acting on $V_\ell A$.  Hence the trace
and degree of $\alpha$ are the trace and determinant of this action.
\end{proposition}

\begin{proof}[Proof of Proposition~\ref{prop:milne-endomorphism-tate-characteristic}]
After replacing $k$ by a separable closure, for every $\beta\in\End(A)$,
\[
|\deg(\beta)|_\ell=|\#\ker(\beta)|_\ell
 =|\operatorname{coker}(T_\ell\beta)|_\ell^{-1}
 =|\det(T_\ell\beta)|_\ell.
\]
Apply this to $\beta=F(\alpha)$ for $F\in\Z[X]$.  If $a_i$ are the roots of
$P_\alpha$ and $b_i$ the eigenvalues on $V_\ell A$, the displayed equality
gives $|\prod_iF(a_i)|_\ell=|\prod_iF(b_i)|_\ell$ for every $F$.  Lemma~
\ref{lem:milne-ell-adic-root-separation} implies that the multisets of roots
agree, hence the characteristic polynomials agree.
\end{proof}

\begin{lemma}[Separation of $\ell$-adic roots]
\label{lem:milne-ell-adic-root-separation}

\dcref{MI:ch1:10.21}
Let $P(X)=\prod_i(X-a_i)$ and $Q(X)=\prod_i(X-b_i)$ be monic polynomials
of the same degree with coefficients in $\Q_\ell$.  If
\[
\left|\prod_iF(a_i)\right|_\ell
=\left|\prod_iF(b_i)\right|_\ell
\qquad(F\in\Z[T]),
\]
then $P=Q$.
\end{lemma}

\begin{proof}
By continuity the equality holds for $F\in\Q_\ell[T]$.  Let $a$ be a root
of $P$ and let $d$ and $e$ be its multiplicities in $P$ and $Q$.  For
$z\in\Q_\ell^{\rm al}$ close to $a$, apply the hypothesis to the minimal
polynomial of $z$ over $\Q_\ell$.  The two sides are respectively a fixed
nonzero factor times $|z-a|_\ell^{md}$ and a fixed nonzero factor times
$|z-a|_\ell^{me}$, where $m$ is the number of distinct conjugates of $z$.
Letting $z\to a$ gives $d=e$.  This holds for every root, so $P=Q$.
\end{proof}

\begin{lemma}[Multiplicative polynomial functions]
\label{lem:milne-multiplicative-polynomial}

\dcref{MI:ch1:10.22}
Let $E$ be an algebra over a field $K$, and let $\delta:E\to K$ be a
polynomial function satisfying $\delta(\alpha\beta)=\delta(\alpha)\delta(\beta)$.
For $\alpha\in E$, let $P(X)=\prod_i(X-a_i)$ be the polynomial such that
$P(x)=\delta(\alpha-x)$.  Then, for every $F\in K[T]$,
\[
\delta(F(\alpha))=\prod_iF(a_i).
\]
\end{lemma}

\begin{proof}
After extending $K$, factor $F(T)=\prod_j(T-b_j)$.  Multiplicativity and
the definition of $P$ give
\[
\delta(F(\alpha))=\prod_jP(b_j)=\prod_{i,j}(b_j-a_i)
 =\prod_iF(a_i).
\]
\end{proof}

\begin{theorem}[Tate-module realization of homomorphisms]
\label{thm:milne-hom-tate-injective}

\dcref{MI:ch1:10.15}
\uses{def:milne-tate-module}
For $\ell\ne\operatorname{char}(k)$, the map
\[
 \Hom(A,B)\otimes\Z_\ell\longrightarrow
 \Hom_{\Z_\ell}(T_\ell A,T_\ell B)
\]
is injective with torsion-free cokernel.  In particular,
$\Hom(A,B)$ is a free abelian group of finite rank at most
$4\dim(A)\dim(B)$.
\end{theorem}

\begin{proof}[Proof of Theorem~\ref{thm:milne-hom-tate-injective}]
We first prove the assertion under the assumption that $\Hom(A,B)$ is
finitely generated over $\Z$.  Let $e_1,\ldots,e_m$ be a basis for
$\Hom(A,B)$, and suppose that
\[
\sum_i a_iT_\ell(e_i)=0,\qquad a_i\in\Z_\ell.
\]
For each $i$, choose a sequence of integers $n_i(r)$ converging
$\ell$-adically to $a_i$.  Then $|n_i(r)|_\ell$ is constant for $r$ large,
i.e., the power of $\ell$ dividing $n_i(r)$ does not change after a certain
point.  But for $r$ large
\[
T_\ell\left(\sum_i n_i(r)e_i\right)=
\sum_i n_i(r)T_\ell(e_i)
\]
is close to zero in $\Hom(T_\ell A,T_\ell B)$, which means that it is
divisible by a high power of $\ell$.  By Lemma~\ref{lem:milne-tate-divisibility},
each $n_i(r)$ is then divisible by a high power of $\ell$.  This contradicts
the preceding assertion unless all $a_i$ are zero.

It remains to prove that $\Hom(A,B)$ is finitely generated over $\Z$.  We
first show that $\End(A)$ is finitely generated when $A$ is simple.  Let
$e_1,\ldots,e_m$ be linearly independent over $\Z$ in $\End(A)$, and let
$P$ be the polynomial function on $\End^0(A)$ such that $P(\alpha)=\deg(\alpha)$.
Because $A$ is simple, a nonzero endomorphism $\alpha$ of $A$ is an isogeny,
and so $P(\alpha)$ is an integer $>0$.  Let $M$ be the $\Z$-submodule of
$\End^0(A)$ generated by the $e_i$.  The map $P:\Q M\to\Q$ is continuous
for the real topology because it is a polynomial in the coordinates.  Thus
\[
U=\{v:P(v)<1\}
\]
is an open neighbourhood of $0$.  Since
$(\Q M\cap\End(A))\cap U=0$, the group $\Q M\cap\End(A)$ is discrete in
$\Q M$, and therefore is a finitely generated $\Z$-module.  Taking the $e_i$
to be a $\Q$-basis for $\End^0(A)$ gives that $\End(A)$ is finitely generated.

For arbitrary $A,B$, choose isogenies
\[
\prod_i A_i^{r_i}\longrightarrow A,\qquad
B\longrightarrow\prod_j B_j^{s_j}
\]
with the $A_i$ and $B_j$ simple.  Then
\[
\Hom(A,B)\longrightarrow\prod_{i,j}\Hom(A_i,B_j)
\]
is injective.  The factors are zero when $A_i$ and $B_j$ are not isogenous;
when they are isogenous, $\Hom(A_i,B_j)$ embeds in a finitely generated
endomorphism group.  This proves finite generation.  Lemma~\ref{lem:milne-tate-divisibility}
gives the torsion-free cokernel, and the rank bound follows from the ranks of
the Tate modules.
\end{proof}

\begin{lemma}[Divisibility detected on torsion]
\label{lem:milne-tate-divisibility}

\dcref{MI:ch1:10.16}
\uses{def:milne-tate-module}
Let $\alpha\in\Hom(A,B)$.  If the homomorphism $T_\ell\alpha$ is divisible
by $\ell^n$ in $\Hom_{\Z_\ell}(T_\ell A,T_\ell B)$, then $\alpha$ is divisible by
$\ell^n$ in $\Hom(A,B)$.
\end{lemma}

\begin{proof}
The hypothesis says that $\alpha$ is zero on $A_{\ell^n}$.  The quotient
property of multiplication by $\ell^n$ therefore gives a homomorphism
$\beta:A\to B$ with $\alpha=\beta\circ\ell_A^n$.
\end{proof}

\begin{proposition}[Characteristic polynomial for a field of endomorphisms]
\label{prop:milne-endomorphism-field-characteristic}

\dcref{MI:ch1:10.23}
\uses{prop:milne-endomorphism-tate-characteristic}
Let $K$ be a $\Q$-subalgebra of $\End(A)\otimes\Q$ with the same identity
element, and assume that $K$ is a field.  If $f=[K:\Q]$, then $V_\ell(A)$ is
a free $K\otimes_\Q\Q_\ell$-module of rank $2\dim(A)/f$.  Consequently, for
$\alpha\in K$,
\[
\Tr(\alpha)=\frac{2g}{f}\Tr_{K/\Q}(\alpha),
\qquad
\deg(\alpha)=\Nm_{K/\Q}(\alpha)^{2g/f}.
\]
\end{proposition}

\begin{proof}[Source proof]
We prove the stronger assertion in which $K$ is replaced by a division
algebra $D$.  Let $K$ be the centre of $D$, put
$d=[D:K]^{1/2}$ and $f=[K:\Q]$.  If $D\otimes_\Q\Q_\ell$ is again a division
algebra, then
\[
V_\ell(A)\simeq V^{\,2\dim(A)/(f d)}
\]
where $V$ is any simple $D\otimes_\Q\Q_\ell$-module.  In general,
$D\otimes_\Q\Q_\ell$ decomposes as a product $\prod_iD_i$ of simple
algebras.  If $D_i\simeq M_{r_i}(D_i^0)$ and $V_i$ is a simple
$M_{r_i}(D_i^0)$-module, then
\[
V_\ell(A)\simeq\bigoplus_i V_i^{m_i}
\]
for some $m_i\geq0$.  We show that the $m_i$ are all equal.

For $\alpha\in D$, let $P_{D/\Q,\alpha}(X)$ be its reduced
characteristic polynomial over $\Q$.  The characteristic polynomial
$P_\alpha(X)$ of $\alpha$ as an endomorphism of $A$ is the characteristic
polynomial of $V_\ell(\alpha)$, and the decomposition gives
\[
P_\alpha(X)=\prod_iP_{D_i/\Q_\ell,\alpha}(X)^{m_i},
\qquad
P_{D/\Q,\alpha}(X)=\prod_iP_{D_i/\Q_\ell,\alpha}(X).
\]
Choose $\alpha$ generating a maximal subfield of $D$, so that
$P_{D/\Q,\alpha}(X)$ is irreducible.  Every monic irreducible factor of
$P_\alpha(X)$ then shares a root with $P_{D/\Q,\alpha}(X)$ and hence equals
it.  Thus $P_\alpha(X)=P_{D/\Q,\alpha}(X)^m$ for an integer $m$, and each
$m_i$ equals $m$.  In the commutative case $D=K$, comparison of dimensions
gives the stated rank, and comparison of trace and norm gives
\[
\Tr(\alpha)=\frac{2g}{f}\Tr_{K/\Q}(\alpha),\qquad
\deg(\alpha)=\Nm_{K/\Q}(\alpha)^{2g/f}.
\]
\end{proof}

\begin{corollary}[Roots of the characteristic polynomial]
\label{cor:milne-characteristic-polynomial-roots}

\dcref{MI:ch1:10.24}
\uses{prop:milne-endomorphism-field-characteristic}
Let $\alpha\in\End^0(A)$ and assume that $\Q[\alpha]$ is a product of fields.
If $C_\alpha(X)$ is the characteristic polynomial of $\alpha$ on $\Q[\alpha]$,
then
\[
\{a\in\C:C_\alpha(a)=0\}=\{a\in\C:P_\alpha(a)=0\}.
\]
\end{corollary}

\begin{proof}
Apply Proposition~\ref{prop:milne-endomorphism-field-characteristic} to
each field factor of $\Q[\alpha]$ and compare the resulting characteristic
polynomials with $P_\alpha$.
\end{proof}

\begin{corollary}[Finite Neron--Severi group]
\label{cor:milne-ns-finite-rank}

\dcref{MI:ch1:10.18}
\uses{thm:milne-endomorphism-characteristic-polynomial}
The Neron--Severi group $\NS(A)$ is a torsion-free finitely generated abelian
group, with
$\operatorname{rank}\NS(A)\leq 4(\dim A)^2$.
\end{corollary}

\begin{proof}
The map $L\mapsto\lambda_L$ injects $\NS(A)$ into
$\Hom(A,A^\vee)$, since its kernel is precisely $\Pic^0(A)$.  Apply the
Tate-module realization theorem.
\end{proof}


\section*{Auxiliary numbered statements (I.5--I.10)}

\begin{lemma}[Base change for sections]
\label{lem:milne-5-8}

\dcref{MI:ch1:5.8}
\lean{MilneLib.LinearMap.bijective_of_surjective_rank_one}
Let $\alpha:M\to N$ be a homomorphism of free modules of rank one over a
commutative ring $R$.  If $\alpha$ is surjective, then it is bijective.
Consequently, a surjective homomorphism of invertible sheaves is an
isomorphism.
\end{lemma}

\begin{proof}

Choose bases $e$ and $f$ and write $\alpha(e)=rf$.  Surjectivity implies
$r\in R^\times$, so multiplication by $r$ is bijective.  Apply this on
stalks for invertible sheaves.

\end{proof}


\begin{lemma}[Affine module and sheaf base change]
\label{lem:milne-5-9}

\dcref{MI:ch1:5.9}
\lean{MilneLib.affineModuleSheaf, MilneLib.affineModuleGlobalSectionsIso,
MilneLib.affineModuleGlobalSectionsIso_naturality, MilneLib.pushforward_globalSections_top,
MilneLib.affineVectorSpaceSheaf, MilneLib.affineVectorSpaceGlobalSectionsIso}
Let $V$ be a variety over $k$ and let $\alpha:V\to\operatorname{Specm}k$ be
the structure map.  A coherent sheaf on $\operatorname{Specm}k$ is a finite
dimensional $k$-vector space $M$, and $\alpha_*\mathcal M=\Gamma(V,\mathcal
M)$ while $\alpha^*M=\mathcal O_V\otimes_k M$.
\end{lemma}

\begin{proof}

This is the definition of pushforward and pullback for the structure map to
the one-point affine scheme.

\end{proof}


\begin{lemma}[Tensor and internal Hom]
\label{lem:milne-5-10}

\dcref{MI:ch1:5.10}
\lean{MilneLib.tensorProductEval, MilneLib.tensorProductEval_tmul,
MilneLib.tensorProductEval_eq_of_tmul, MilneLib.schemeSheafEvaluation,
MilneLib.schemeSheafEvaluationAt, MilneLib.schemeSheafEvaluation_app}
\leanok
For a homomorphism $R\to S$ of commutative rings and an $S$-module $M$, the
map $S\otimes_RM\to M$, $s\otimes m\mapsto sm$, is canonical.  Likewise,
for a regular map $\alpha:W\to V$ and a coherent $\mathcal O_W$-module
$\mathcal M$, there is a canonical evaluation map
$\alpha^*\alpha_*\mathcal M\to\mathcal M$.
\end{lemma}

\begin{proof}
\leanok
On an affine open this is the usual multiplication map
$\mathcal O_W\otimes_R M\to M$.
\end{proof}

\begin{lemma}[Exactness for locally free sheaves]
\label{lem:milne-5-11}

\dcref{MI:ch1:5.11}
\lean{MilneLib.moduleFinite_residueFieldTensor,
MilneLib.moduleFinite_schemeModuleStalk_residueTensor,
MilneLib.SchemeModule.IsStalkwiseFinite.residueFieldTensor}
\uses{lem:milne-5-8}
Let $\alpha:M\to N$ be a homomorphism of $R$-modules, with $N$ finitely
generated.  If for every maximal ideal $\mathfrak m$ the induced map
$M/\mathfrak mM\to N/\mathfrak mN$ is surjective, then
$M_\mathfrak m\to N_\mathfrak m$ is surjective.  For coherent sheaves, if
the maps on all residue fibres are surjective, then the sheaf map is
surjective; if both sheaves are invertible, it is an isomorphism.
\end{lemma}

\begin{proof}

This is Nakayama's lemma, applied after localization; the sheaf statement is
the same argument on stalks, followed by Lemma~\ref{lem:milne-5-8}.

\end{proof}


\begin{lemma}[Cohomology after field extension]
\label{lem:milne-5-12}

\dcref{MI:ch1:5.12}
Let $V$ be a complete variety over $k$ and let $\mathcal M$ be a locally free
sheaf.  For every field extension $K/k$, if $\mathcal M'$ is its pullback to
$V_K$, then
\[
\Gamma(V_K,\mathcal M')=\Gamma(V,\mathcal M)\otimes_kK.
\]
For a divisor $D$ on a smooth complete variety,
$L(D_K)=L(D)\otimes_kK$.
\end{lemma}

\begin{proof}

This is the flat base-change statement for global sections on a complete
variety (equivalently, Theorem~\ref{thm:milne-cohomology-base-change} in
degree zero).

\end{proof}


\begin{lemma}[Relative fibre notation]
\label{lem:milne-5-13}

\dcref{MI:ch1:5.13}
Let $V$ be complete and let $\mathcal L$ be locally free on $V$.  If
$\mathcal L$ becomes trivial after a field extension $K/k$, then it is
already trivial on $V$.
\end{lemma}

\begin{proof}
An invertible sheaf on a complete variety is trivial exactly when it and its
dual have nonzero global sections.  Apply Lemma~\ref{lem:milne-5-12} to
both sheaves and descend the nonzero sections.
\end{proof}

\begin{lemma}[Fibrewise constancy criterion]
\label{lem:milne-5-14}

\dcref{MI:ch1:5.14}
Consider a regular map $V\to T$.  For $t\in T$, write
$V_t=V\times_k\operatorname{Specm}k(t)$ and write $\mathcal L_t$ for the
pullback of a sheaf $\mathcal L$ on $V\times T$ to $V_t$.
\end{lemma}

\begin{proof}

This is the usual fibre notation; when $k$ is algebraically closed,
$k(t)=k$ and the fibres are literal copies of $V$.

\end{proof}


\begin{lemma}[Seesaw reduction]
\label{lem:milne-5-15}

\dcref{MI:ch1:5.15}
Let $\alpha:V\to T$ be proper.  Then $\alpha_*\mathcal M$ is coherent for
every coherent sheaf $\mathcal M$ on $V$.  If $\mathcal L$ is invertible on
$V\times T$ with $V$ complete, then $t\mapsto\dim_{k(t)}\Gamma(V_t,\mathcal
L_t)$ is upper semicontinuous.  If it is constantly $n$, then $q_*\mathcal L$
is locally free of rank $n$ and the canonical map
$(q_*\mathcal L)(t)\to\Gamma(V_t,\mathcal L_t)$ is an isomorphism.
\end{lemma}

\begin{proof}
The coherence and semicontinuity assertions are the proper base-change theorem;
Milne refers to Mumford for the proof.
\end{proof}
\begin{proposition}[Linear systems separate points]
\label{prop:milne-6-1}

\dcref{MI:ch1:6.1}
Let $\mathfrak d$ be a complete linear system on a complete nonsingular variety
$V$, and let $\phi:V\dashrightarrow\PP^n$ be the rational map defined by
$\mathfrak d$.  Then $\phi$ is defined at $P$ if and only if $P$ is not a base
point of $\mathfrak d$.
\end{proposition}

\begin{proof}
Choose $D_0\in\mathfrak d$ and a basis $f_0,\ldots,f_n$ of $L(D_0)$.  If
$P$ is not a base point, choose $D_0$ with $P\notin\operatorname{Supp}(D_0)$.
Since the system has no fixed divisor, each $f_i/f_0$ has no pole at $P$ and
the projective point $[f_0(P):\cdots:f_n(P)]$ is defined.  Conversely, if $P$
is a base point, every divisor in the system contains $P$, so all homogeneous
coordinates vanish at $P$ and the rational map is undefined there.
\end{proof}

\begin{definition}[Separation by a linear system]
\label{def:milne-6-2}

\dcref{MI:ch1:6.2}
\begin{enumerate}
\item A linear system $\mathfrak d$ separates points if, for every pair of
points $P,Q\in V$, there is a $D\in\mathfrak d$ with
$P\in\operatorname{Supp}(D)$ and $Q\notin\operatorname{Supp}(D)$.
\item It separates tangent directions if, for every $P\in V$ and every
nonzero tangent vector $t$ at $P$, there is a divisor $D\in\mathfrak d$ with
$P\in D$ but $t\notin T_P(D)$.
\end{enumerate}
\end{definition}

\begin{proposition}[Criterion for an embedding]
\label{prop:milne-6-3}

\dcref{MI:ch1:6.3}
\uses{prop:milne-6-1,def:milne-6-2}
Assume that $\mathfrak d$ has no base points.  Then the map
$\phi:V\to\PP^n$ defined by $\mathfrak d$ is a closed immersion if and only if
$\mathfrak d$ separates points and separates tangent directions.
\end{proposition}

\begin{proof}[Source proof omitted]
The necessity is immediate from the preceding discussion.  For sufficiency,
Milne refers to Hartshorne 1977, II 7.8; no proof is included in the notes.
\end{proof}


\begin{proposition}[Ample divisors on an abelian variety]
\label{prop:milne-6-6}

\dcref{MI:ch1:6.6}
\uses{prop:milne-6-3,thm:milne-square}
For divisors on a complete nonsingular variety:
\begin{enumerate}
\item if $D$ and $D'$ are ample, then $D+D'$ is ample;
\item if $D$ is ample, then its restriction to a closed subvariety is ample
whenever defined;
\item $D$ is ample exactly when its scalar extension to $k^{\rm al}$ is ample;
\item if $V_{k^{\rm al}}$ has an ample divisor, then $V$ has an ample divisor.
\end{enumerate}
\end{proposition}

\begin{proof}
For (a), choose $n$ such that both $nD$ and $nD'$ are very ample.  Since
$nD'$ is very ample, it is linearly equivalent to an effective divisor $D''$.
Then $L(nD+D'')\supset L(nD)$, so $nD+D''$ is very ample and hence so is
$nD+nD'$.  For (b), restrict the map defined by $D$ to $W$; this is the map
defined by $D|_W$.  For (c), scalar extension of the map $V\to\PP^n$ defined
by $D$ is the map defined by $D_{k^{\rm al}}$ (Lemma~\ref{lem:milne-5-12}).
For (d), let $D$ be ample on $V_{k^{\rm al}}$.  It is defined over a finite
extension $k'/k$, so the set of its distinct Galois conjugates is finite.  By
(a), their sum $D'$ is ample.  It is defined over a finite purely inseparable
extension of $k$; if $k$ is imperfect, a suitable Frobenius multiple of $D'$
is defined over $k$.
\end{proof}


\begin{lemma}[Dual correspondence of a line bundle]
\label{lem:milne-8-8}

\dcref{MI:ch1:8.8}
\uses{def:milne-dual-abelian,thm:milne-square}
For an invertible sheaf $\mathcal L$ on $A$, the class
$[t_a^*\mathcal L\otimes\mathcal L^{-1}]$ depends algebraically on $a$
and defines a homomorphism $\lambda_{\mathcal L}:A\to A^\vee$.
\end{lemma}

\begin{proof}

The theorem of the square gives the homomorphism identity; the normalized
Poincare bundle identifies the resulting family with a regular map to the dual.

\end{proof}

\begin{proposition}[Coherent sheaves on a finite quotient]
\label{prop:milne-8-13}

\dcref{MI:ch1:8.13}
\uses{prop:milne-finite-quotient}
Assume that the finite group $G$ acts freely on $V$, and let $W=V/G$.  Then
$\mathcal M\mapsto\pi_*\mathcal M$ defines an equivalence from the category
of coherent $\mathcal O_W$-modules to the category of coherent $G$-sheaves on
$V$, under which locally free sheaves of rank $r$ correspond to locally free
sheaves of rank $r$.
\end{proposition}

\begin{proof}[Source proof omitted]
Milne refers to Mumford, p. 70, for the descent equivalence for coherent
sheaves under a free finite quotient.
\end{proof}

% Source-faithful blueprint fragment, Milne PDF pp. 53--83 (PDF 59--89).
% Pedagogical notes, exercises, and applications are omitted; statements retain
% the printed numbering and source/dcref metadata.

\section{Polarizations and invertible sheaves}
\begin{definition}[Polarization]\label{def:milne-polarization}\dcref{MI:ch1:11}
A polarization $\lambda$ of an abelian variety $A$ is an isogeny $\lambda:A\to A^\vee$ which over $k^{\rm al}$ is $\lambda_L$ for an ample sheaf $L$ on $A_{k^{\rm al}}$. Its degree is its degree as an isogeny. A polarized abelian variety is a pair $(A,\lambda)$; degree one gives a principal polarization.
\end{definition}

\begin{theorem}[Riemann--Roch for abelian varieties]\label{thm:milne-polarization-riemann-roch}\dcref{MI:ch1:11.1}
Let $L$ be an invertible sheaf on $A$, and put $\chi(L)=\sum_i(-1)^i\dim_kH^i(A,L)$. Then (a) $\deg(\lambda_L)=\chi(L)^2$; (b) if $L=L(D)$, then $\chi(L)=(D^g)/g!$; (c) if $\dim K(L)=0$, then $H^r(A,L)$ is nonzero for exactly one integer $r$.
\end{theorem}

\begin{proof}[Source proof omitted]
Milne states this as Mumford's Riemann--Roch theorem and does not prove it in
these notes; he refers to Mumford for the proof.
\end{proof}

\begin{theorem}[Finiteness over finite fields]\label{thm:milne-polarization-finiteness}\dcref{MI:ch1:11.2}
If $k$ is finite and $g,d>0$, then up to isomorphism there are only finitely many abelian varieties $A/k$ of dimension $g$ possessing a polarization of degree $d^2$.
\end{theorem}

\begin{proof}

Theorem~11.1 realizes $A$ as a variety of degree $3^gd(g!)$ in $\PP^{3^gd-1}$. Its Chow form is homogeneous of that degree in each of $g+1$ sets of variables and determines the isomorphism class. There are only finitely many such forms over $k$.

\end{proof}


\section{Etale cohomology of an abelian variety}
\begin{theorem}[Etale cohomology]\label{thm:milne-etale-cohomology}\dcref{MI:ch1:12.1}
Let $A$ have dimension $g$ over a separably closed field $k$, and let $\ell\ne\operatorname{char}(k)$. There is a canonical isomorphism $H^1(A_{\rm et},\Z_\ell)\simeq\operatorname{Hom}_{\Z_\ell}(T_\ell A,\Z_\ell)$, and cup products give $\bigwedge^rH^1(A_{\rm et},\Z_\ell)\simeq H^r(A_{\rm et},\Z_\ell)$ for every $r$. In particular $H^r$ is free of rank $\binom{2g}{r}$.
\end{theorem}
\begin{proof}[Proof of Theorem~\ref{thm:milne-etale-cohomology}]
When $R=\Q_\ell$ or $\F_\ell$, the conditions of the Hopf algebra lemma are
fulfilled with $d=2g$ (Milne 1980, VI 1.1).  Therefore
$\dim H^1(A,\Q_\ell)\leq2g$.  The canonical injection
\[
\Hom_{\Z_\ell}(T_\ell A,\Z_\ell)\hookrightarrow H^1(A,\Z_\ell)
\]
and the fact that $T_\ell A$ has rank $2g$ show that equality holds.  The
lemma now gives
\[
H^r(A,\Q_\ell)\simeq\bigwedge^rH^1(A,\Q_\ell),
\]
and in particular its dimension is $\binom{2g}{r}$.  Applying the same
argument with $\F_\ell$ and using the exact coefficient sequence
\[
\cdots\longrightarrow H^r(A,\Z_\ell)\xrightarrow{\ell}
H^r(A,\Z_\ell)\longrightarrow H^r(A,\F_\ell)
\longrightarrow H^{r+1}(A,\Z_\ell)\longrightarrow\cdots
\]
shows that $H^r(A,\Z_\ell)$ is torsion-free for every $r$.  Consequently the
map from $\bigwedge^rH^1(A,\Z_\ell)$ to $H^r(A,\Z_\ell)$ is injective after
tensoring with $\Q_\ell$ and surjective after reducing modulo $\ell$, hence is
an isomorphism.
\end{proof}
\begin{lemma}[Hopf algebra lemma]\label{lem:milne-etale-hopf}\dcref{MI:ch1:12.2}
Let $H^*$ be a graded associative anticommutative algebra over a perfect field $K$, with $H^0=K$, vanishing above degree $d$, and a map $m^*$ satisfying $m^*(x)=x\otimes1+1\otimes x+\sum x_i\otimes y_i$ with positive degrees in the sum. Then $\dim H^1\le d$; equality implies that $H^*$ is canonically the exterior algebra on $H^1$.
\end{lemma}

\begin{proof}
Borel's structure theorem for Hopf algebras shows that $H^*$ is the
associative algebra generated by certain elements $x_i$, subject only to the
relations imposed by anticommutativity and nilpotence.  The product of the
$x_i$ has degree $\sum_i\deg(x_i)$, and therefore $\sum_i\deg(x_i)\leq d$.
In particular, the number of generators of degree one is at most $d$; this
number is $\dim H^1$.  If equality holds, every $x_i$ has degree one.  Their
squares must be zero, since otherwise a nonzero product with one repeated
factor would have degree $d+1$.  Thus $H^*$ is canonically the exterior
algebra on $H^1$.
\end{proof}


\section{Weil pairings}
\begin{lemma}[Compatibility of Weil pairings]\label{lem:milne-weil-pairing-compatibility}\dcref{MI:ch1:13.1}
If $m,n$ are prime to $\operatorname{char}(k)$, then $e_{mn}(a,a')^n=e_m(na,na')$ for $a\in A_{mn}(k)$ and $a'\in A^\vee_{mn}(k)$.
\end{lemma}

\begin{proof}[Source proof omitted]
Milne refers to Milne 1986, 16.1 (or leaves this as an exercise), and does not
include the proof in these notes.
\end{proof}
\begin{proposition}[Functoriality of Weil pairings]\label{prop:milne-weil-pairing-functoriality}\dcref{MI:ch1:13.2}
For $\alpha:A\to B$, $e_\ell(a,\alpha^\vee b)=e_\ell(\alpha a,b)$; $e_\ell^{\alpha^\vee\lambda\alpha}(a,a')=e_\ell^\lambda(\alpha a,\alpha a')$; and $e_\ell^{\alpha^*L}(a,a')=e_\ell^L(\alpha a,\alpha a')$. Moreover $L\mapsto e_\ell^L$ is a homomorphism $\Pic(A)\to\operatorname{Hom}(\bigwedge^2T_\ell A,\Z_\ell(1))$, so $e_\ell^L$ is skew-symmetric.
\end{proposition}

\begin{proof}[Source proof omitted]
Milne refers to Milne 1986, 16.2 (or leaves this as an exercise), and does not
include the proof in these notes.
\end{proof}
\begin{theorem}[Descent of Neron--Severi classes]\label{thm:milne-weil-pairing-descent}\dcref{MI:ch1:13.4}
If $\alpha:A\to B$ is an isogeny of degree prime to $\operatorname{char}(k)$ and $\lambda\in\NS(A)$, then $\lambda=\alpha^*(\lambda')$ for some $\lambda'\in\NS(B)$ iff for every $\ell\mid\deg(\alpha)$ there is a skew-symmetric form $e$ on $T_\ell B$ with $e_\ell^\lambda(a,a')=e(\alpha a,\alpha a')$.
\end{theorem}

\begin{proof}[Source proof omitted]
Milne lists this result and refers to Milne 1986, 16.4 (or Mumford, sections
20 and 23) for its proof.
\end{proof}
\begin{corollary}[Divisibility criterion]\label{cor:milne-weil-pairing-divisibility}\dcref{MI:ch1:13.5}
For $\ell\ne\operatorname{char}(k)$, $\lambda\in\NS(A)$ is divisible by $\ell^n$ iff $e_\ell^\lambda$ is divisible by $\ell^n$ in $\operatorname{Hom}(\bigwedge^2T_\ell A,\Z_\ell(1))$.
\end{corollary}

\begin{proof}
Apply Theorem~\ref{thm:milne-weil-pairing-descent} to $\ell^n_A$.
\end{proof}

\begin{proposition}[Polarization from skew pairing]\label{prop:milne-weil-pairing-skew}\dcref{MI:ch1:13.6}
If $\operatorname{char}(k)\ne2$, a homomorphism $\lambda:A\to A^\vee$ is $\lambda_L$ for some $L\in\Pic(A)$ iff $e_\ell^\lambda$ is skew-symmetric.
\end{proposition}
\begin{proof}[Source proof omitted]
Milne lists this result and refers to the Weil-pairing construction for the
proof.
\end{proof}
\begin{lemma}[Poincare pairing]\label{lem:milne-poincare-pairing}\dcref{MI:ch1:13.7}
For the Poincare sheaf $P$ on $A\times A^\vee$, $e_\ell^P((a,b),(a',b'))=e_\ell(a,b')e_\ell(a',b)^{-1}$.
\end{lemma}
\begin{proof}[Source proof omitted]
Milne states this formula as part of the Weil-pairing discussion and does not
include a proof.
\end{proof}
\begin{proposition}[Descent of a polarization]\label{prop:milne-polarization-descent}\dcref{MI:ch1:13.8}
Let $\alpha:A\to B$ be an isogeny of degree prime to $\operatorname{char}(k)$ and $\lambda$ a polarization of $A$. Then $\lambda=\alpha^*\lambda'$ for a polarization $\lambda'$ on $B$ iff $\ker(\alpha)\subset\ker(\lambda)$ and $e^\lambda$ is trivial on $\ker(\alpha)\times\ker(\alpha)$.
\end{proposition}
\begin{proof}[Source proof omitted]
Milne lists this criterion and refers to the preceding descent theorem and the
standard polarization argument for its proof.
\end{proof}
\begin{corollary}[Principal polarization up to isogeny]\label{cor:milne-principal-polarization}\dcref{MI:ch1:13.10}
An abelian variety with a polarization of degree prime to the characteristic is isogenous to a principally polarized abelian variety.
\end{corollary}

\begin{proof}

Choose a prime $\ell$ dividing the degree and a subgroup $N\subset\ker(\lambda)$ of order $\ell$. Skew-symmetry makes $e^\lambda$ trivial on $N\times N$; Proposition~\ref{prop:milne-polarization-descent} lowers the degree by $\ell^2$. Iterate.

\end{proof}

\begin{corollary}[Zarhin auxiliary polarization]\label{cor:milne-zarhin-auxiliary}\dcref{MI:ch1:13.11}
If $\ker(\lambda)\subset A_m$ with $m$ prime to the characteristic and $\alpha\in\End(A)$ preserves $\ker(\lambda)$ and satisfies $\alpha^\vee\lambda\alpha=\lambda$ on $A_{m^2}$, then $A\times A^\vee$ is principally polarized.
\end{corollary}
\begin{proof}[Source proof omitted]
Milne states this auxiliary result without proof and uses it in the proof of
Zarhin's trick.
\end{proof}
\begin{theorem}[Zarhin's trick]\label{thm:milne-zarhin-trick}\dcref{MI:ch1:13.12}
For every abelian variety $A$, $(A\times A^\vee)^4$ is principally polarized.
\end{theorem}

\begin{proof}
Let $\lambda$ be a polarization of $A$, and let $m$ kill $\ker(\lambda)$.
By Lagrange's four-square theorem, choose integers $a,b,c,d$ such that
$a^2+b^2+c^2+d^2\equiv1\pmod {m^2}$.  Set
\[
\alpha=\begin{pmatrix}
a&b&c&d\\ b&a&d&c\\ c&d&a&b\\ d&c&b&a
\end{pmatrix}\in M_4(\Z)\subset\End(A^4).
\]
It commutes with $\lambda^4=\operatorname{diag}(\lambda,\lambda,lambda,\lambda)$,
so $\alpha(\ker(\lambda^4))\subset\ker(\lambda^4)$.  Moreover, $\alpha^\vee$
is the transpose of $\alpha$, and hence
\[
\alpha^\vee\circ\lambda^4\circ\alpha
=\lambda^4\circ\alpha^{\mathsf t}\alpha
=(a^2+b^2+c^2+d^2)\lambda^4
\]
on $A_{m^2}$.  The auxiliary polarization result therefore applies to
$A^4$ and proves the assertion.
\end{proof}

\begin{corollary}[Finite-dimensional abelian varieties over finite fields]\label{cor:milne-finite-abelian-varieties}\dcref{MI:ch1:13.13}
Over a finite field, for each $g$ there are only finitely many isomorphism classes of abelian varieties of dimension $g$.
\end{corollary}

\begin{proof}

Apply Theorem~\ref{thm:milne-zarhin-trick} and the degree-one case of Theorem~\ref{thm:milne-polarization-finiteness} in dimension $8g$; then use the finiteness of direct factors.

\end{proof}


\section{The Rosati involution}
\begin{definition}[Rosati involution]\label{def:milne-rosati-involution}\dcref{MI:ch1:14}
For a polarization $\lambda$ on $A$, put $\alpha^\dagger=\lambda^{-1}\circ\alpha^\vee\circ\lambda$ on $\End(A)\otimes\Q$.
\end{definition}
\begin{proposition}[Neron--Severi classes and Rosati fixed points]\label{prop:milne-rosati-ns}\dcref{MI:ch1:14.2}
Over an algebraically closed field, $L\mapsto\lambda^{-1}\lambda_L$ identifies $\NS(A)\otimes\Q$ with the elements of $\End(A)\otimes\Q$ fixed by $\dagger$.
\end{proposition}

\begin{proof}

By Proposition~13.6, $\lambda\alpha=\lambda_L$ iff $e_\ell^\lambda(a,\alpha a')=e_\ell^\lambda(\alpha a',a)$, equivalently $\lambda\alpha=\alpha^\vee\lambda$, i.e. $\alpha=\alpha^\dagger$.

\end{proof}

\begin{theorem}[Rosati positivity]\label{thm:milne-rosati-positivity}\dcref{MI:ch1:14.3}
The form $(\alpha,\beta)\mapsto\Tr(\alpha\beta^\dagger)$ on $\End(A)\otimes\Q$ is positive definite. If $D$ defines $\lambda$, then $\Tr(\alpha\alpha^\dagger)=\frac{2g}{(D^g)}(D^{g-1}\cdot\alpha^*D)>0$ for $\alpha\ne0$.
\end{theorem}
\begin{proof}[Source proof omitted]
Milne proves positivity by the displayed intersection calculation but says that
the calculation is omitted here (see Milne 1986, section 17).
\end{proof}
\begin{proposition}[Automorphisms of polarized varieties]\label{prop:milne-rosati-automorphisms}\dcref{MI:ch1:14.4}
The automorphism group of $(A,\lambda)$ is finite. For $n\ge3$, an automorphism acting identically on $A_n(k^{\rm al})$ is the identity.
\end{proposition}

\begin{proof}
For an automorphism $\alpha$ of $(A,\lambda)$ we have
$\lambda=\alpha^\vee\lambda\alpha$, and hence $\alpha^\dagger\alpha=1$.
Consequently
\[
\operatorname{Aut}(A,\lambda)\subset
\End(A)\cap\{x\in\End(A)\otimes\R:\Tr(x^\dagger x)=2g\}.
\]
The second set is compact by Rosati positivity, while $\End(A)$ is discrete;
this proves (a).  For (b), if $\alpha$ acts identically on $A_n$, then
$\alpha-1=n\beta$ for some $\beta\in\End(A)$ by Lemma
\ref{lem:milne-tate-divisibility}.  The eigenvalues of $\alpha$ and $\beta$
are algebraic integers, and those of $\alpha$ are roots of unity because
$\alpha$ has finite order.  The preceding lemma shows that all eigenvalues of
$\alpha$ equal $1$.  Thus $\alpha-1=n\beta$ is nilpotent.  If $\beta\ne0$,
put $\beta' =\beta^\dagger\beta\ne0$; positivity gives
$\Tr(\beta^\dagger\beta)>0$.  Since $\beta'$ is fixed by the Rosati
involution, $\Tr((\beta')^2)>0$, and inductively $(\beta')^{2^r}\ne0$ for
every $r$, contradicting nilpotence.  Hence $\beta=0$ and $\alpha=1$.
\end{proof}


\section{Geometric finiteness theorems}
\begin{lemma}[Root of unity lemma]\label{lem:milne-root-unity}\dcref{MI:ch1:14.5}
If $\zeta$ is a root of unity and $\zeta=1+n\alpha$ for an algebraic integer $\alpha$ and integer $n\ge3$, then $\zeta=1$.
\end{lemma}

\begin{proof}

Let $p$ be a prime divisor of the order of $\zeta$ and replace $\zeta$ by a
power of prime order $p$.  The norm of $\zeta-1$ is $p$ up to sign, whereas
$\zeta-1=n\alpha$ makes this norm divisible by $n^{p-1}$; for $n\ge3$ this
is impossible.  Thus the order has no prime divisor and $\zeta=1$.

\end{proof}

\begin{theorem}[Finitely many polarizations]\label{thm:milne-finite-polarizations}\dcref{MI:ch1:15.1}
Let $A$ be an abelian variety over a field $k$, and let $d$ be an integer.
Then there are only finitely many isomorphism classes of polarized abelian
varieties $(A,\lambda)$ with $\deg\lambda=d$.
\end{theorem}
\begin{proof}[Source proof]
Fix a polarization $\lambda_0$ of $A$ and let $\dagger$ be the Rosati
involution defined by $\lambda_0$.  The map
\[
\lambda\longmapsto\lambda_0^{-1}\circ\lambda
\]
identifies the polarizations of $A$ with a subset of the $\dagger$-fixed
elements of $\End(A)\otimes\Q$.  Choose an isogeny $\alpha:A^\vee\to A$
with $\alpha\circ\lambda_0=n$ for some $n\in\Z$.  Then
$\lambda_0^{-1}=(n_A^{-1})\circ\alpha$, so the image lies in the lattice
$L=n^{-1}\End(A)$.  The action of $\End(A)^\times$ on polarizations is
\[
\lambda\longmapsto u^\vee\circ\lambda\circ u,
\]
and under the preceding map it becomes $x\mapsto u^\dagger xu$.
Moreover, $\deg(\lambda_0^{-1}\lambda)=\deg(\lambda_0)^{-1}\deg(\lambda)$;
the characteristic-polynomial and reduced-norm calculation in Section~I.10,
this gives a bound on the norm of the corresponding elements of $L$ when
$\deg(\lambda)$ is fixed.  Proposition~\ref{prop:milne-arithmetic-orbits} therefore gives
finitely many orbits, and Proposition~\ref{prop:milne-rosati-automorphisms}
turns these into finitely many isomorphism classes.
\end{proof}
\begin{proposition}[Arithmetic orbit finiteness]\label{prop:milne-arithmetic-orbits}\dcref{MI:ch1:15.2}
Let $E$ be a finite-dimensional semisimple $\Q$-algebra with involution, $R$ an order, and $L$ an $R^\times$-stable lattice under $x\mapsto u^\dagger xu$. For each $N$, the set $\{v\in L:\Norm(v)\le N\}/R^\times$ is finite.
\end{proposition}
\begin{proof}[Source proof omitted]
Milne says that this proposition is proved using the reduction theory of
arithmetic subgroups and gives the applications below; no proof is included.
\end{proof}
\begin{theorem}[Finitely many direct factors]\label{thm:milne-finite-direct-factors}\dcref{MI:ch1:15.3}
Up to isomorphism, an abelian variety has only finitely many direct factors.
\end{theorem}

\begin{proof}
Let $B$ be a direct factor of $A$, with complement $C$.  Under an isomorphism
$A\simeq B\times C$, define $e$ by the composite
\[
A\simeq B\times C\longrightarrow B\times C,
\qquad (b,c)\longmapsto(b,0),
\]
followed by the inverse isomorphism.  Then $e$ is an idempotent and $B$ is
determined by $e$, since $B=\ker(1-e)$.  Conversely, for every idempotent
$e\in\End(A)$,
\[
A=\ker(1-e)\times\ker(e).
\]
If $u\in\End(A)^\times$, then $e'=ueu^{-1}$ is an idempotent and $u$ gives
an isomorphism $\ker(1-e)\to\ker(1-e')$.  Thus there is a surjection
\[
\{\text{idempotents in }\End(A)\}/\End(A)^\times
\twoheadrightarrow
\{\text{direct factors of }A\}/\simeq.
\]
The assertion follows from Proposition~\ref{prop:milne-idempotent-orbits}.
\end{proof}

\begin{proposition}[Integral idempotent orbits]\label{prop:milne-idempotent-orbits}\dcref{MI:ch1:15.4}
If $E$ is semisimple finite-dimensional over $\Q$ and $R$ an order, then the idempotents in $R$ form finitely many $R^\times$-conjugacy classes.
\end{proposition}
\begin{proof}[Source proof omitted]
Milne says that this proposition is proved by applying the general reduction
theorem below (see his application 15.8); no proof is included in the notes.
\end{proof}
\begin{theorem}[Borel--Harish-Chandra orbit theorem]\label{thm:milne-bhc-orbit}\dcref{MI:ch1:15.5}
Let $G/\Q$ be reductive, $\Gamma\subset G(\Q)$ arithmetic, $G\to\GL(V)$ a representation, and $L$ a $\Gamma$-stable lattice. If $X$ is a closed $G$-orbit in $V$, then $L\cap X$ is a finite union of $\Gamma$-orbits.
\end{theorem}
\begin{proof}[Source proof omitted]
Milne cites Borel 1969, 9.11 for this theorem and does not include its proof.
\end{proof}
\section{Families of abelian varieties}
\begin{proposition}[Rigidity lemma]\label{prop:milne-rigidity}\dcref{MI:ch1:16.1}
Let $S$ be a connected variety, let $\pi:V\to S$ be proper and flat with
geometrically connected fibres, let $\pi':V'\to S$ be a variety over $S$,
and let $\alpha:V\to V'$ be a morphism over $S$.  If the image of one fibre
$V_s$ in $V'_s$ is a single point, then $\alpha$ factors through $S$: there is
a map $\alpha':S\to V'$ with $\alpha=\alpha'\circ\pi$.
\end{proposition}
\begin{proof}[Source proof omitted]
Mumford, D., *Geometric Invariant Theory*, Springer, 1965, 6.1.
\end{proof}
\begin{corollary}[Families are group-homomorphic]\label{cor:milne-family-homomorphism}\dcref{MI:ch1:16.2}
\begin{enumerate}
\item Every morphism of families of abelian varieties carrying the zero
section into the zero section is a homomorphism.
\item The group structure on a family of abelian varieties is uniquely
determined by the choice of a zero section.
\item A family of abelian varieties is commutative.
\end{enumerate}
\end{corollary}
\begin{proof}
For (a), apply the rigidity lemma to the map
$A\times_SA\to B$ defined by the difference between the two sides of the
homomorphism identity.  Part (b) follows immediately.  For (c), the inverse
map $a\mapsto-a$ is a homomorphism by (a).
\end{proof}
\begin{proposition}[Constant-family subvarieties]\label{prop:milne-constant-family-subvariety}\dcref{MI:ch1:16.3}
Let $A$ be an abelian variety over a field $k$, and let $S$ be a variety over
$k$ such that $S(k)\ne\varnothing$.  For any injective homomorphism of
families $\alpha:B\hookrightarrow A\times S$ over $S$, there is an abelian
subvariety $B_0$ of $A$, defined over $k$, such that
$\alpha(B)=B_0\times S$.
\end{proposition}
\begin{proof}
Let $s\in S(k)$, and let $B_0=B_s$ be the fibre over $s$.  Then $\alpha_s$
identifies $B_0$ with a subvariety of $A$.  The map
\[
h:B\hookrightarrow A\times S\longrightarrow (A/B_0)\times S
\]
has fibre $B_0\to A\to A/B_0$ over $s$, which is zero, and so the Rigidity
Lemma shows that $h=0$.  It follows that $\alpha(B)\subset B_0\times S$.
In fact, the two are equal because their fibres over $S$ are connected and
have the same dimension.
\end{proof}
\begin{corollary}[Regular extensions]\label{cor:milne-regular-extension}\dcref{MI:ch1:16.4}
Let $K$ be a regular extension of a field $k$.
\begin{enumerate}
\item Let $A$ be an abelian variety over $k$.  Then every abelian subvariety
of $A_K$ is defined over $k$.
\item If $A$ and $B$ are abelian varieties over $k$, then every homomorphism
$A_K\to B_K$ is defined over $k$.
\end{enumerate}
\end{corollary}
\begin{proof}
Let $V$ be a variety over $k$ such that $k(V)=K$.  After replacing $V$ by an
open subvariety, we may assume that the abelian subvariety extends to a family
of abelian varieties over $V$.  If $V$ has a $k$-rational point, the preceding
proposition shows that it is defined over $k$.  In any case, there is a finite
Galois extension $k'/k$ and an abelian subvariety $B'$ of $A_{k'}$ such that
\[
B_{Kk'}=B'_{Kk'}.
\]
This equality uniquely determines $B'$ as a subvariety of $A_{k'}$.  Since
$\sigma B'$ has the same property for every $\sigma\in\operatorname{Gal}(k'/k)$,
we have $\sigma B'=B'$, and hence $B'$ descends to $k$.  For (b), part (a)
shows that the graph of $\alpha$ is defined over $k$.
\end{proof}
\begin{theorem}[$K/k$-trace]\label{thm:milne-trace}\dcref{MI:ch1:16.5}
Let $K/k$ be a regular extension of fields and let $A$ be an abelian variety
over $K$.  Then there exist an abelian variety $B$ over $k$ and a
homomorphism $\alpha:B_K\to A$ with finite kernel such that, for every
abelian variety $B'$ over $k$ and every homomorphism
$\alpha':B'_K\to A$ with finite kernel, there is a unique homomorphism
$\varphi:B'\to B$ with $\alpha'=\alpha\circ\varphi_K$.
\end{theorem}

\begin{proof}
Consider the collection of pairs $(B,\alpha)$ in the statement, and let
$A^*$ be the abelian subvariety of $A$ generated by the images of the
$\alpha$.  For two pairs $(B_1,\alpha_1)$ and $(B_2,\alpha_2)$, the identity
component $C$ of the kernel of
\[
(\alpha_1,\alpha_2):(B_1\times B_2)_K\longrightarrow A
\]
is an abelian subvariety of $B_1\times B_2$, and Corollary
\ref{cor:milne-regular-extension} shows that it is defined over $k$.  The map
$(B_1\times B_2)/C\to A$ has finite kernel and image the subvariety generated
by $\alpha_1(B_1)$ and $\alpha_2(B_2)$.  It is now clear that there is a pair
$(B,\alpha)$ whose image is $A^*$.  Divide $B$ by the largest subgroup scheme
$N$ of $\ker(\alpha)$ defined over $k$.  The resulting pair $(B/N,\alpha)$
has the required universal property, as in Milne's construction.

\end{proof}

\begin{theorem}[Mordell--Weil]\label{thm:milne-mordell-weil}\dcref{MI:ch1:16.7}
If $K$ is finitely generated over the prime field and $A/K$ is abelian, then $A(K)$ is finitely generated.
\end{theorem}
\begin{proof}[Source proof omitted]
Milne gives the historical attribution (Mordell, Weil, Taniyama, Lang--Neron)
and refers to Serre 1989, Chapter 4, for a proof over number fields; no proof
of the general theorem is included here.
\end{proof}
\begin{theorem}[Lang--Neron]\label{thm:milne-lang-neron}\dcref{MI:ch1:16.8}
If $K/k$ is finitely generated with $k$ algebraically closed in $K$, and $(B,\alpha)$ is the $K/k$-trace of $A$, then $A(K)/\alpha(B)(k)$ is finitely generated.
\end{theorem}
\begin{proof}[Source proof omitted]
Milne states this theorem and refers to Conrad 2006, Hindry (Appendix B to
Kahn 2006), and Kahn 2006 for proofs; no proof is included in the reference.
\end{proof}

\section{Neron models and semistable reduction}
\begin{theorem}[Neron model special fibre]\label{thm:milne-neron-model}\dcref{MI:ch1:17.1}
Let $R$ be a discrete valuation ring with fraction field $K$ and residue field
$k$, and let $A$ be an abelian variety over $K$.  There is a canonical way
to attach to $A$ an algebraic group $A^0$ over $k$.
\end{theorem}
\begin{proof}[Source proof omitted]
Milne states the theorem and records the representing Neron model in a
subsequent remark, but does not prove representability here.
\end{proof}
\begin{theorem}[Semistable reduction]\label{thm:milne-semistable-reduction}\dcref{MI:ch1:17.3}
If $A$ has good reduction, then $A^0$ does not change under a finite field
extension.  If $A$ has semistable reduction, then $(A^0)^0$ does not change
under a finite field extension.  Every abelian variety acquires semistable
reduction after a finite extension.
\end{theorem}
\begin{proof}[Source proof]
These are the Neron--Ogg--Shafarevich and semistable reduction theorems; their
proofs are long and difficult, and Milne gives the statements here without a
proof.
\end{proof}

\section{Abel and Jacobi}
\begin{theorem}[Holomorphic differentials]\label{thm:milne-holomorphic-differentials}\dcref{MI:ch1:18.1}
For a compact Riemann surface $C$ of genus $g$, the holomorphic differentials form a $g$-dimensional complex vector space.
\end{theorem}
\begin{proof}[Source proof omitted]
Milne states this standard theorem without proof.
\end{proof}
\begin{theorem}[Independence of periods]\label{thm:milne-period-independence}\dcref{MI:ch1:18.2}
Let $\omega_1,\ldots,\omega_g$ be a basis of the holomorphic differentials
on a compact Riemann surface $C$, and let $\gamma_1,\ldots,\gamma_{2g}$ be
the standard basis of $H_1(C,\Z)$.  The period vectors
\[
\pi_j=\left(\int_{\gamma_j}\omega_1,\ldots,
\int_{\gamma_j}\omega_g\right)\in\C^g
\]
are linearly independent over $\R$.
\end{theorem}
\begin{proof}[Source proof omitted]
Milne states the independence of the period vectors without proof.
\end{proof}
\begin{theorem}[Jacobi inversion]\label{thm:milne-jacobi-inversion}\dcref{MI:ch1:18.3}
Let $\ell:\Omega(C)\to\C$ be a linear map.  Then there exist points
$P_1,\ldots,P_g$ on $C$ and paths $\gamma_i$ from a fixed point $P$ to
$P_i$ such that
\[
\ell(\omega)=\sum_{i=1}^g\int_{\gamma_i}\omega
\qquad\text{for all }\omega\in\Omega(C).
\]
\end{theorem}
\begin{proof}[Source proof omitted]
Milne states Jacobi inversion without proof.
\end{proof}
\begin{theorem}[Abel's theorem]\label{thm:milne-abel-theorem}\dcref{MI:ch1:18.4}
Let $P_1,\ldots,P_r$ and $Q_1,\ldots,Q_r$ be points of $C$, not
necessarily distinct.  There is a meromorphic function on $C$ with poles
exactly at the $P_i$ and zeros exactly at the $Q_i$ if and only if, for every
choice of paths $\gamma_i$ from $P$ to $P_i$ and $\gamma_i'$ from $P$ to
$Q_i$, there is a $\gamma\in H_1(C,\Z)$ such that
\[
\sum_i\int_{\gamma_i}\omega-\sum_i\int_{\gamma_i'}\omega
 =\int_{\gamma}\omega
\qquad\text{for all }\omega\in\Omega(C).
\]
\end{theorem}
\begin{proof}[Source proof omitted]
Milne states Abel's theorem without proof.
\end{proof}
\begin{theorem}[Jacobian from periods]\label{thm:milne-period-jacobian}\dcref{MI:ch1:18.5}
Put
\[
J=\Omega(C)^\vee/H_1(C,\Z),
\]
where $H_1(C,\Z)$ is embedded by integration.  The intersection product
\[
H_1(C,\Z)\times H_1(C,\Z)\longrightarrow\Z
\]
is a Riemann form on $J$; hence $J$ is an abelian variety.  Fixing $P\in C$,
the map
\[
\sum_i n_iP_i\longmapsto\left(\omega\longmapsto
\sum_i n_i\int_P^{P_i}\omega\right):\Div^0(C)\longrightarrow J
\]
is surjective with kernel the principal divisors, and therefore induces an
isomorphism $\Pic^0(C)\simeq J$.
\end{theorem}
\begin{proof}
The intersection pairing is integral, alternating, and positive for the
complex structure by the Riemann bilinear relations, so it is a Riemann form.
Jacobi inversion gives surjectivity of the divisor map, and Abel's theorem
identifies its kernel with principal divisors.  Passing to the quotient gives
the asserted Picard isomorphism.
\end{proof}
